% coord_R2_oblic.tex
%
% Copyright (C) 2022--2026 José A. Navarro Ramón <janr.devel@gmail.com>
% Licencia del código GPLv2
% Licencia Creative Commons Recognition Non-Commercial Share-alike.
% (CC-BY-NC-SA)

\chapter{Sistema de coordenadas cartesianas oblicuas en
  \mathinhead{\symbb{R}^2}{espr2}}
Este capítulo tiene los siguientes objetivos:
\begin{enumerate}
\item Analizar el sistema de coordenadas oblicuas en el plano euclídeo $\symbb{R}^2$ y
  determinadas operaciones, utilizando unos conceptos básicos de geometría diferencial.
\item Realizar algunas operaciones diferenciales, como gradiente, divergencia y
  laplaciano en este sistema.
\end{enumerate}

Se utilizarán determinados conceptos de la geometría diferencial, como \emph{base natural}, \emph{métrica}, \emph{factores de escala},
\emph{vectores y componentes contravariantes y covariantes},
\emph{espacio dual}, etc., conforme se van necesitando, unos con más detalle que otros,
de manera que, al final del capítulo, se termina con una idea básica de esta rama de las
matemáticas. Por descontado, la geometría diferencial es mucho más que lo que aquí se
presenta, pues solo se ha arañado una pequeña parte de la superficie.

%Aquí introduciremos los conceptos de \emph{componentes contravariantes y covariantes},
%algo que no hicimos en el sistema cartesiano rectangular, porque en este último ambas
%componentes son idénticas, pero en este otro sí hay diferencia.
%Solo adelantaremos que las componentes contravariantes se representan situando el índice
%de la coordenada como superíndice (no es un exponente), y la covariantes poniéndolo
%como subíndice. Más adelante se aclararán estos términos.

\section{Visión general. Características fundamentales del sistema}
\begin{itemize}
\item Ejes y su notación:
  Tendremos dos ejes; uno de ellos, horizontal, lo haremos coincidir con el eje $x$
  del sistema de coordenadas ortogonales, mientras que el otro formará un ángulo $\alpha$
  con el primero. Es como si el eje $y$ del sistema ortogonal hubiera girado en sentido
  horario un ángulo $\beta = \qty{90}-\alpha$ si $\alpha$ es menor de \qty{90}{\degree}, o
  $\beta = \alpha - \qty{90}{\degree}$ en caso contrario.
  
  En $\symbb{R}^2$ utilizaremos indistintamente, dependiendo del contexto, dos
  notaciones diferentes para los ejes, ver figura \ref{fig:R2-oblic_ejes}:
  \begin{itemize}
  \item Notación en álgebra lineal y física clásica:
    El eje horizontal se suele representar como $\tilde{x}$ o $x'$ y el inclinado como
    $\tilde{y}$ o $y'$.
    Esta notación se suele usar para diferenciar visualmente dos sistemas cartesianos,
    como el ortogonal estándar y el oblicuo, o entre dos oblicuos.
    En este texto, a menos que se diga lo contrario, cuando utilicemos esta notación,
    reservaremos $x$ e $y$ para el sistema cartesiano ortogonal, aunque el sistema
    oblicuo no lo comparemos con aquél.
  \item Notación en cálculo tensorial y geometría diferencial:
    Se usan preferentemente superíndices $x^1$ y $x^2$
    (que no representan exponentes, ni componentes contravariantes de vectores, que
    también se pueden representar igual, aunque se distinguen por el contexto).
  \end{itemize}
\item Base natural:
  La más utilizada en geometría diferencial es la base covariante, formada por vectores
  que se representan mediante subíndices
  \[
    B = \set{\xhat{e}_1, \xhat{e}_2}
  \]
  Los vectores se han representado con acento circunflejo porque en este sistema son
  vectores unitarios, aunque esto no es cierto en otros sistemas de coordenadas.

  Avanzamos un detalle importante: los vectores covariantes $\xhat{e}_i$ (con subíndice)
  de la base son los que sirven para medir las componentes contravariantes $x^i$ (con
  superíndice) de los vectores, como se aprecia en la figura \ref{fig:R2-oblic_contrav}.
\end{itemize}
\begin{figure}[ht]
  % Escala
  \def\scl{1.00}
  %  Longitud de los ejes
  \pgfmathsetmacro{\EJECARTLONG}{3.0}
  \pgfmathsetmacro{\EJEOBLICLONG}{2.5}
  % Eje x
  \pgfmathsetmacro{\XANG}{0}
  \pgfmathsetmacro{\XLONG}{\EJECARTLONG}
  % Eje y
  \pgfmathsetmacro{\YANG}{90}
  \pgfmathsetmacro{\YLONG}{\EJECARTLONG}
  % Eje xtilde
  \pgfmathsetmacro{\XTILDEANG}{0}
  \pgfmathsetmacro{\XTILDELONG}{\EJEOBLICLONG}    
  % Eje ytilde
  \pgfmathsetmacro{\YTILDEANG}{70}
  \pgfmathsetmacro{\YTILDELONG}{0.5*(\EJECARTLONG+\EJEOBLICLONG)}
  % Fondo
  \pgfmathsetmacro{\HORZ}{0.90}
  \pgfmathsetmacro{\VERT}{0.25}
  % 
  \centering
  \begin{tikzpicture}[%
    scale=\scl,
    every node/.style={black,font=\small},
    ejecartesiano/.style={->, black!30},
    ejeoblicuo/.style={->, black},
    textooblicuo/.style={black},
    textocartesiano/.style={\colorTorig},
    angulo/.style={fill=green!15, draw=\colorGirodos},
    angulogiro/.style={draw=black!50},
    background/.style={%
      line width=\bgborderwidth,
      draw=\bgbordercolor,
      fill=\bgcolor,
    },
    ]
    % COORDENADAS
    % Origen
    \coordinate (O) at (0,0);
    % Eje x
    \path[save path=\ejex] (O) -- ++(\XANG:\XLONG cm) coordinate (xfin);
    % Eje y
    \path[save path=\ejey] (O) -- ++(\YANG:\YLONG cm) coordinate (yfin);
    % Eje xtilde
    \path[save path=\ejextilde] (O) -- ++(\XTILDEANG:\XTILDELONG cm)
    coordinate (xtildefin);
    % Eje ytilde
    \path[save path=\ejeytilde, name path=O--ytilde] (O) -- ++(\YTILDEANG:\YTILDELONG cm)
    coordinate (ytildefin);

    % DIBUJOS
    % Ángulos
    % Ángulo alfa
    \path (xfin) -- (O) -- (ytildefin) pic
    [angulo, "\footnotesize $\alpha$",angle radius=14pt, angle eccentricity=0.64]
    {angle = xfin--O--ytildefin};
    % Ángulo beta
    \path (yfin) -- (O) -- (ytildefin) pic
    [angulogiro, {Latex[round,width=3pt,length=4pt]}-,
    "\color{black!60}{\footnotesize $\beta$}", angle radius=50pt,
    angle eccentricity=0.80]
    {angle = ytildefin--O--yfin};
    
    % Ejes
    % Eje x
    \draw[ejecartesiano, use path=\ejex];
    \node[textocartesiano, below] at (xfin) {$x$};
    % Eje y
    \draw[ejecartesiano, use path=\ejey];
    \node[textocartesiano, left] at (yfin) {$y$};
    
    % Eje xtilde
    \draw[ejeoblicuo, use path=\ejextilde];
    \node[below, textooblicuo] at (xtildefin) {$\tilde{x}$};
    % Eje ytilde
    \draw[ejeoblicuo, use path=\ejeytilde];
    %\node[above right=0pt and -4pt, texto] at (ytildefin) {$\tilde{x}$};
    \node[above right=0pt and -4pt, textooblicuo] at (ytildefin) {$\tilde{y}$};

    % Origen
    \fill[black] (O) circle [radius=0.8pt];
    
    % Fondo amarillo
    \coordinate (SW) at ($(current bounding box.south west) + (-\HORZ cm,-\VERT cm)$);
    \coordinate (NE) at ($(current bounding box.north east) + (\HORZ cm,\VERT cm)$);
    \begin{scope}[on background layer]
      \draw[background] (SW) rectangle (NE);
    \end{scope}
  \end{tikzpicture}
  % Escala
  \def\scl{1.00}
  %  Longitud de los ejes
  \pgfmathsetmacro{\EJECARTLONG}{3.0}
  \pgfmathsetmacro{\EJEOBLICLONG}{2.5}
  % Eje x
  \pgfmathsetmacro{\XANG}{0}
  \pgfmathsetmacro{\XLONG}{\EJECARTLONG}
  % Eje y
  \pgfmathsetmacro{\YANG}{90}
  \pgfmathsetmacro{\YLONG}{\EJECARTLONG}
  % Eje xtilde
  \pgfmathsetmacro{\XTILDEANG}{0}
  \pgfmathsetmacro{\XTILDELONG}{\EJEOBLICLONG}    
  % Eje ytilde
  \pgfmathsetmacro{\YTILDEANG}{70}
  \pgfmathsetmacro{\YTILDELONG}{0.5*(\EJECARTLONG+\EJEOBLICLONG)}
  % Fondo
  \pgfmathsetmacro{\HORZ}{0.90}
  \pgfmathsetmacro{\VERT}{0.25}
  % 
  \centering
  \begin{tikzpicture}[%
    scale=\scl,
    every node/.style={black,font=\small},
    ejecartesiano/.style={->, black!30},
    ejeoblicuo/.style={->, black},
    textooblicuo/.style={black},
    textocartesiano/.style={\colorTorig},
    angulo/.style={fill=green!15, draw=\colorGirodos},
    angulogiro/.style={draw=black!50},
    background/.style={%
      line width=\bgborderwidth,
      draw=\bgbordercolor,
      fill=\bgcolor,
    },
    ]
    % COORDENADAS
    % Origen
    \coordinate (O) at (0,0);
    % Eje x
    \path[save path=\ejex] (O) -- ++(\XANG:\XLONG cm) coordinate (xfin);
    % Eje y
    \path[save path=\ejey] (O) -- ++(\YANG:\YLONG cm) coordinate (yfin);
    % Eje x^1
    \path[save path=\ejextilde] (O) -- ++(\XTILDEANG:\XTILDELONG cm)
    coordinate (xtildefin);
    % Eje x^2
    \path[save path=\ejeytilde, name path=O--ytilde] (O) -- ++(\YTILDEANG:\YTILDELONG cm)
    coordinate (ytildefin);
    
    % DIBUJOS
    % Ángulo alpha
    \path (xfin) -- (O) -- (ytildefin) pic
    [angulo, "\footnotesize $\alpha$",angle radius=14pt, angle eccentricity=0.64]
    {angle = xfin--O--ytildefin};
    % Ángulo beta
    \path (yfin) -- (O) -- (ytildefin) pic
    [angulogiro, {Latex[round,width=3pt,length=4pt]}-,
    "\color{black!60}{\footnotesize $\beta$}", angle radius=50pt,
    angle eccentricity=0.80]
    {angle = ytildefin--O--yfin};
    
    % Ejes
    % Eje x
    \draw[ejecartesiano, use path=\ejex];
    \node[textocartesiano, below] at (xfin) {$x$};
    % Eje y
    \draw[ejecartesiano, use path=\ejey];
    \node[textocartesiano, left] at (yfin) {$y$};
    
    % Eje x^1
    \draw[ejeoblicuo, use path=\ejextilde];
    \node[below, textooblicuo] at (xtildefin) {$x^1$};
    % Eje x^2
    \draw[ejeoblicuo, use path=\ejeytilde];
    \node[above right=0pt and -4pt, textooblicuo] at (ytildefin) {$x^2$};

    % Origen
    \fill[black] (O) circle [radius=0.8pt];
    
    % Fondo amarillo
    \coordinate (SW) at ($(current bounding box.south west) + (-\HORZ cm,-\VERT cm)$);
    \coordinate (NE) at ($(current bounding box.north east) + (\HORZ cm,\VERT cm)$);
    \begin{scope}[on background layer]
      \draw[background] (SW) rectangle (NE);
    \end{scope}
  \end{tikzpicture}
  \caption{A la izquierda los ejes oblicuos se marcan como $\tilde{x}$ y $\tilde{y}$,
    mientras que a la derecha se encuentran como $x^1$ y $x^2$.
  Estos forman un ángulo $\alpha$ entre sí. Los ejes ortogonales estándar $x$ e $y$ se
  representan para comparar. Nótese que se hace coincidir el eje $x^1$ con el $x$.
  Por último, se puede considerar que el eje $y$ se ha inclinado un ángulo en sentido
  horario $\beta = \pi/2 - \alpha$.}
  \label{fig:R2-oblic_ejes}
\end{figure}

\subsection{Base natural}
Como se comentó en el apartado anterior, la base natural, ahora $B'$ en el sistema
oblicuo (con prima, para distinguirla temporalmente de la base ortogonal
$B=\set{\uvec{\i},\,\uvec{\j}}$ con la que la compararemos más adelante), está formada
por los vectores covariantes (nótese el subíndice) siguientes
\begin{equation}\label{eq:R2-oblic_base}
  B' = \set{\xhat{e}_1, \xhat{e}_2}
\end{equation}
Esta base es muy importante, pues servirá para obtener la métrica de estos sistemas
oblicuos que, a su vez, nos permitirá relacionar de alguna manera las componentes
covariantes y las contravariantes, entre sí.

\subsubsection{Cálculo geométrico}
Hasta que no obtengamos la matriz de cambio de base más adelante, para obtener la base
natural que necesitamos, vamos a recurrir a un procedimiento geométrico. Para ello,
empezamos fijándonos en la figura \ref{fig:R2-coord_rect-ij}.
De ella deducimos que, en coordenadas ortogonales, la posición de cualquier punto del
plano euclídeo se expresa como
\begin{equation}
  \vvv{r} = x\,\uvec{\i} + y\,\uvec{\j}
\end{equation}

La base natural se calcula expresando el vector de posición, expresado en función de las
componentes contravariantes en el nuevo sistema $x^1$ y $x^2$, y calculando las
derivadas parciales $\vvv{e}_i = \partial\vvv{r}/\partial x^i$.
Hasta que no se explique con algo de detalle el significado de estas componentes
contravariantes, es suficiente con fijarnos en la figura \ref{fig:R2-coord_rect-ij}.
De acuerdo con ella, el vector de posición de un punto cualquiera del plano se expresa
como
\[
  \vvv{r} = x\,\uvec{\i} + y\,\uvec{\j}
\]

Nuestro objetivo es obtener $x$ e $y$ en función de las $x^1$ y $x^2$
\begin{align}
  x &= x(x^1, x^2)\\
  y &= y(x^1, x^2)
\end{align}
y sustituir en la expresión anterior, para poder derivar parcialmente el vector de
posición respecto de las $x^i$.

Fijémonos en el triángulo sombreado en verde de la figura \ref{fig:R2-oblic_proced_geom}.
Sabemos que el ángulo $\beta$ es el complementario de $\alpha$, que es el que forman los
dos ejes oblicuos, de manera que $\sin\beta = \cos\alpha$ y $\cos\beta = \sin\alpha$.

Así
\begin{equation}
  \begin{cases}
    \begin{array}{lll}
    \sin\beta = \cos\alpha = \dfrac{x-x^1}{x^2}
      &\longrightarrow
      &x = x^1 + x^2\cos\alpha\\
      \cos\beta = \sin\alpha = \dfrac{y}{x^2}
      &\longrightarrow
      &y = x^2\sin\alpha
    \end{array}
  \end{cases}
\end{equation}

\begin{figure}[ht]
  % Escala
  \def\scl{1.00}
  % 
  \pgfmathsetmacro{\EJECARTLONG}{3.5}
  \pgfmathsetmacro{\EJEOBLICLONG}{3.0}
  % Eje x
  \pgfmathsetmacro{\XANG}{0.0}
  \pgfmathsetmacro{\XLONG}{\EJECARTLONG}
  % Eje y
  \pgfmathsetmacro{\YANG}{90}
  \pgfmathsetmacro{\YLONG}{\EJECARTLONG}
  % Eje xtilde
  \pgfmathsetmacro{\XTILDEANG}{0}
  \pgfmathsetmacro{\XTILDELONG}{\EJEOBLICLONG}    
  % Eje ytilde
  \pgfmathsetmacro{\YTILDEANG}{55}
  \pgfmathsetmacro{\YTILDELONG}{0.56*(\EJECARTLONG + \EJEOBLICLONG)}
  % Vector P
  \pgfmathsetmacro{\PMOD}{3.0}
  \pgfmathsetmacro{\PANG}{30}
  % Fondo
  \pgfmathsetmacro{\HORZ}{0.25}
  \pgfmathsetmacro{\VERT}{0.25}
  % 
  \centering
 \begin{tikzpicture}[%
    scale=\scl,
    every node/.style={black,font=\footnotesize},
    ejecart/.style={->, black!30},
    ejeoblic/.style={->, black},
    proyeccionoblic/.style={black!20, line width=0.5},
    proyeccioncart/.style={black!20, line width=0.5, dashed},
    % longitud/.style={line width=.5pt,\colorV},
    texto/.style={black},
    textoapagado/.style={\colorTorig},
    vector/.style={-{Latex[round]}, shorten >=1.2pt, line width=.8pt,\colorV},
    contravariante/.style={line width=1.4pt, red},
    cartesiana/.style={gray, line width=1.5pt},
    contravtxt/.style={red!70!black},
    pcirculo/.style={fill=\colorPorig, draw=black},
    angulo/.style={fill=green!15, draw=\colorGirodos},
    angulogiro/.style={draw=black!50},
    background/.style={%
      line width=\bgborderwidth,
      draw=\bgbordercolor,
      fill=\bgcolor,
    },
    ]
    % COORDENADAS
    % Origen
    \coordinate (O) at (0,0);
    % Eje x
    \path[save path=\ejex] (O) -- ++(\XANG:\XLONG cm) coordinate (xfin);
    % Eje x
    \path[save path=\ejey] (O) -- ++(\YANG:\YLONG cm) coordinate (yfin);
    % Eje xtilde
    \path[save path=\ejextilde] (O) -- ++(\XTILDEANG:\XTILDELONG cm)
    coordinate (xtildefin);
    % Eje ytilde
    \path[save path=\ejeytilde, name path=O--ytilde] (O) -- ++(\YTILDEANG:\YTILDELONG cm)
    coordinate (ytildefin);
    % Vector
    \path[save path=\r] (O) -- (\PANG:\PMOD cm) coordinate (P);
    % Proyección de P en los ejes
    % Con eje x
    \path[save path=\proyx] (P) |- (xfin) coordinate[pos=0.5] (px);
    % Con eje y
    \path[save path=\proyy] (P) -| (yfin) coordinate[pos=0.5] (py);
    % Componentes contravariantes
    % Intersección de horizontal desde P hasta ytilde
    % -> Componente contravariante eje ytilde (CTV2) 
    \path[name path=P--y] (P) -| (yfin);
    \path[name intersections={of=P--y and O--ytilde, by=CTV2}];
    % Paralela (CTV2) -- (O) que pase por (P)
    % -> Componente contravariante eje xtilde (CTV1)
    \path (P) -- ($(P) - (CTV2) - (O)$) coordinate (CTV1);
    % Componente delta x = x-x^1
    \path[save path=\deltaxsegmento] (CTV1) -- coordinate[pos=0.5] (deltaxtxt) (px);
    % Componente y
    \path (P) -- coordinate[pos=0.5] (ytxt) (px);

    % DIBUJOS    
    % Área
    \fill[green!15] (P) -- (px) -- (CTV1) -- cycle;

    % Ángulo beta
    \path (yfin) -- (O) -- (ytildefin) pic
    [angulogiro, {Latex[round,width=3pt,length=4pt]}-,
    "\color{black!60}{\footnotesize $\beta$}", angle radius=75pt,
    angle eccentricity=1.1]
    {angle = ytildefin--O--yfin};
    % Otro ángulo beta
    \path (px) -- (P) -- (CTV1) pic
    [angulo, "\color{black}{\scriptsize $\beta$}", angle radius=16pt,
    angle eccentricity=0.70]
    {angle = CTV1--P--px};
    
    % Proyecciones de componentes
    % Con el eje x
    \draw[proyeccioncart, use path=\proyx];
    % Con el eje y
    \draw[proyeccioncart, use path=\proyy];
    % Contravariantes
    \draw[proyeccionoblic] (P) -- (CTV1);
    %\node[below] at (CTV1) {\footnotesize $\hat{x}^1$};
    \draw[proyeccionoblic] (P) -- (CTV2);
    % \node[left] at (CTV2) {\footnotesize $\hat{x}^2$};
    % Eje x
    \draw[ejecart, use path=\ejex];
    \node[textoapagado, below=3pt] at (xfin) {$x$};
    % Eje y
    \draw[ejecart, use path=\ejey];
    \node[textoapagado, left] at (yfin) {$y$};
    % Eje xtilde
    \draw[ejeoblic, use path=\ejextilde];
    \node[texto, below] at (xtildefin) {$x^1$};
    % Eje ytilde
    \draw[ejeoblic, use path=\ejeytilde];
    \node[above right=-1pt and -3pt, texto] at (ytildefin) {$x^2$};
    % Vector v
    \draw[vector, use path=\r];
    \node[above right=-2pt and -2pt] at (P) {\footnotesize $\vvv{r}$};
    \filldraw[fill=green, draw=black] (P) circle[radius=1.2pt];
    %\node[above right=0pt and -5pt] at (P) {\footnotesize $P$};

    % Marcar en color la componentes contravariantes
    \draw[contravariante] (O) -- node[contravtxt, pos=0.5, below] {$x^1$} (CTV1);
    \draw[contravariante] (O) -- node[contravtxt, pos=0.6, left] {$x^2$} (CTV2);

    % delta x
    \draw[cartesiana, use path=\deltaxsegmento];
    \node[below] at (deltaxtxt) {$x-x^1$};
    % y
    %\draw[cartesiana, use path=\ysegmento];
    %\node[right] at (ytxt) {$y$};
    \path (px) -- node[pos=0.41, right] {$y$} (P);
    % x^2
    \path (CTV1) -- node[pos=0.43, left] {$x^2$} (P);

    % Origen
    \fill[orange!80!black] (O) circle [radius=0.8pt];
    
    % Fondo amarillo
    \coordinate (SW) at ($(current bounding box.south west) + (-\HORZ cm,-\VERT cm)$);
    \coordinate (NE) at ($(current bounding box.north east) + (\HORZ cm,\VERT cm)$);
    \begin{scope}[on background layer]
      \draw[background] (SW) rectangle (NE);
    \end{scope}
  \end{tikzpicture}
  \caption{El triángulo rectángulo sombreado en verde nos permite relacionar las
    coordenadas ortogonales $x$ e $y$ con las contravariantes $x^1$ y $x^2$, a través
    del ángulo $\beta$ que forma el eje $x^2$ con el eje $y$.}
  \label{fig:R2-oblic_proced_geom}
\end{figure}

El vector de posición de un punto en función de las componentes contravariantes es
\begin{equation}
  \vvv{r} = \left(x^1+x^2\cos\alpha\right)\,\uvec{\i} + x^2\sin\alpha\,\uvec{j}
\end{equation}

Los vectores de la base natural serán
\begin{align}
  \xhat{e}_1
  &=
    \dfrac{\partial\vvv{r}}{x^1}
    = \uvec{\i}
    \longrightarrow
  |\xhat{e}_1| = |\uvec{\i}| = \sqrt{1} = 1\\
  \xhat{e}_2
  &=
    \dfrac{\partial\vvv{r}}{x^2}
    = \cos\alpha\,\uvec{\i} + \sin\alpha\,\uvec{\j}
    \longrightarrow
    |\xhat{e}_2| = \sqrt{\cos^2\alpha + \sin^2\alpha} = 1
\end{align}
y sus módulos valen la unidad.

\section{Cambio de base}
Para relacionar esta base $B'$ con la estándar de las rectilíneas
$B=\set{\uvec{\i},\uvec{\j}}$, haremos coincidir el versor $\xhat{e}_1$ con $\uvec{\i}$,
esto es, el eje $\tilde{x}$ coincide con el eje $x$ de las cartesianas. Véase la figura
\ref{fig:R2-coord_oblicuas_e1e2}.
\begin{figure}[ht]
  \centering
  % Escala
  \def\scl{1.00}
  % 
  % \pgfmathsetmacro{\EJESCARTLONG}{3.0}
  % \pgfmathsetmacro{\EJEOBLICLONG}{2.7}
  \pgfmathsetmacro{\EJESCARTLONG}{2.5}
  \pgfmathsetmacro{\EJEOBLICLONG}{2.1}
  % Eje x
  \pgfmathsetmacro{\XLONG}{\EJESCARTLONG}
  % Eje y
  \pgfmathsetmacro{\YLONG}{\EJESCARTLONG}
  % Eje xtilde
  \pgfmathsetmacro{\XTILDEANG}{0}
  \pgfmathsetmacro{\XTILDELONG}{\EJEOBLICLONG}    
  % Eje ytilde
  \pgfmathsetmacro{\YTILDEANG}{70}
  \pgfmathsetmacro{\YTILDELONG}{\EJESCARTLONG}
  % Versores
  \pgfmathsetmacro{\VTILDEMOD}{1.4}
  \pgfmathsetmacro{\VXTILDEANG}{0}
  \pgfmathsetmacro{\VYTILDEANG}{\YTILDEANG}
  % Vector P
  \pgfmathsetmacro{\PMOD}{2.3}
  \pgfmathsetmacro{\PANG}{40}
  % Fondo
  \pgfmathsetmacro{\HORZ}{0.25}
  \pgfmathsetmacro{\VERT}{0.25}
  % 
  \centering
  \begin{tikzpicture}[%
    scale=\scl,
    every node/.style={black,font=\small},
    ejecart/.style={->, black!30},
    ejeoblic/.style={->, black},
    proyeccion/.style={black!20, line width=0.4},
    % longitud/.style={line width=.5pt,\colorV},
    texto/.style={black,font=\footnotesize},
    textoapagado/.style={\colorTorig},
    vtilde/.style={%
      -{Latex[round,width=4pt,length=6pt]}, line width=1pt, \colorVred},
    pcirculo/.style={fill=\colorPorig, draw=black},
    angulo/.style={fill=green!15, draw=\colorGirodos},
    background/.style={%
      line width=\bgborderwidth,
      draw=\bgbordercolor,
      fill=\bgcolor,
    },
    ]
    % COORDENADAS
    % Origen
    \coordinate (O) at (0,0);
    % Eje x
    \path[save path=\ejex] (O) -- (\XLONG cm, 0) coordinate (xfin);
    % Eje y
    \path[save path=\ejey] (O) -- (0, \YLONG cm) coordinate (yfin);
    % Eje xtilde
    \path[save path=\ejextilde] (O) -- ++(\XTILDEANG:\XTILDELONG cm)
    coordinate (xtildefin);
    % Eje ytilde
    \path[save path=\ejeytilde] (O) -- ++(\YTILDEANG:\YTILDELONG cm)
    coordinate (ytildefin);
    % Versores tilde
    \path[save path=\vxtilde] (O) -- coordinate[pos=1.0] (vxtildetxt)
    ++(\VXTILDEANG:\VTILDEMOD cm);
    \path[save path=\vytilde] (O) -- coordinate[pos=1.1] (vytildetxt)
    ++(\VYTILDEANG:\VTILDEMOD cm) coordinate (e2fin);
    % Proyección de P en los ejes
    % Con eje x
    \path[save path=\proyx] (e2fin) |- (xfin) coordinate[pos=0.5] (e2x) (xfin);
    % Posición de texto en eje x
    \path (O) -- coordinate[pos=0.5] (extxt) (e2x);
    % Con eje y
    \path[save path=\proyy] (e2fin) -| coordinate[pos=0.5] (e2y) (yfin);
    % Posición de texto en eje y
    \path (O) -- coordinate[pos=0.5] (eytxt) (e2y);
    
    % DIBUJOS
    % Ángulo alpha
    \path (xfin) -- (O) -- (ytildefin) pic
    [angulo, "\footnotesize $\alpha$",angle radius=12pt, angle eccentricity=0.68]
    {angle = xfin--O--ytildefin};
    
    % Proyección de etilde2 en los ejes y texto correspondiente en ejes x e y
    \draw[proyeccion, use path=\proyx];
    \node[below] at (extxt) {\footnotesize $\cos\alpha$};
    \draw[proyeccion, use path=\proyy];
    \node[rotate=90, below left=4pt and 11pt] at (e2y) {\footnotesize $\sin\alpha$};
    
    % Ejes
    % x
    \draw[ejecart, use path=\ejex];
    \node[textoapagado, below=1.2pt] at (xfin) {\small $x$};
    % \node[below] at (xtexto) {$r\cos\theta$};
    % y
    \draw[ejecart, use path=\ejey];
    \node[above, textoapagado] at (yfin) {$y$};
    % xtilde
    \draw[ejeoblic, use path=\ejextilde];
    %\node[texto, below] at (xtildefin) {\footnotesize $\tilde{x}^1$};
    \node[texto, below] at (xtildefin) {\footnotesize $\tilde{x}$};
    % ytilde
    \draw[ejeoblic, use path=\ejeytilde];
    %\node[above right=0pt and -4pt, texto] at (ytildefin) {$\tilde{x}^2$};
    \node[above right=0pt and -4pt, texto] at (ytildefin) {$\tilde{y}$};
    % \node[rotate=90, left=8pt] at (ytexto) {$r\sin\theta$};
    % Versores
    \draw[vtilde, use path=\vxtilde];
    \node[\colorTred, below] at (vxtildetxt) {\footnotesize $\xhat{e}_1$};
    \draw[vtilde, use path=\vytilde];
    \node[\colorTred,above left=-4pt and -1.5pt] at (vytildetxt)
    {\footnotesize $\xhat{e}_2$};
    % Origen
    \filldraw (O) circle [radius=.2pt];
    
    % Fondo amarillo
    \coordinate (SW) at ($(current bounding box.south west) + (-\HORZ cm,-\VERT cm)$);
    \coordinate (NE) at ($(current bounding box.north east) + (\HORZ cm,\VERT cm)$);
    \begin{scope}[on background layer]
      \draw[background] (SW) rectangle (NE);
    \end{scope}
  \end{tikzpicture}
  \caption{Sistema de coordenadas cartesianas oblicuas $\tilde{x}$ e $\tilde{y}$, que
    forman un ángulo $\alpha$ entre sí. También se representan los ejes cartesianos $x$ e
    $y$ tradicionales, como referencia en tono algo apagado. Por último se aprecian en
    rojo los vectores unitarios $\xhat{e}_1$ y $\xhat{e}_2$ del sistema oblicuo y sus
    proyecciones en los ejex $x$ e $y$.}
  \label{fig:R2-coord_oblicuas_e1e2}
\end{figure}

Vamos a expresar estos vectores en función de los $\uvec{\i}$ y $\uvec{\j}$ de las
cartesianas rectangulares
\begin{align}
  \label{eq:R2-e1-ij}
  \xhat{e}_1 &= 1\,\uvec{\i} + 0\,\uvec{\j} = \uvec{\i}\\
  \label{eq:R2-e2-ij}
  \xhat{e}_2 &= \cos{\alpha}\, \uvec{\i} + \sin{\alpha}\, \uvec{\j}
\end{align}
Nótese que los vectores de la base son unitarios, por esta razón los habíamos
representado, ya desde el principio, con el acento circunflejo
\begin{align}
  |\xhat{e}_1|
  &=
    |\uvec{\i}| = 1\\
  |\xhat{e}_2|
  &=
    (\cos\theta\,\uvec{\i}+\sin\theta\,\uvec{\j})^2
    = \cos^2\theta + \sin^2\theta = 1
\end{align}

La matriz que permite obtener la nueva base $B'$ a partir de la antigua $B$, está formada
por las componentes de los vectores $(1,0)$ y $(\cos\alpha,\sin\alpha)$ obtenidos en las
expresiones \eqref{eq:R2-e1-ij} y \eqref{eq:R2-e2-ij}, dispuestas en columnas
\[
  \mmm{M}
  =
  \begin{pmatrix}
    1 & \cos\alpha\\
    0 & \sin\alpha
  \end{pmatrix}
\]

Utilizamos esta matriz para obtener los vectores de la base oblícua a partir de la
rectangular. Es importante darse cuenta que los vectores se representan en filas
\[
  \begin{pmatrix}
    \xhat{e}_1 &  \xhat{e}_2
  \end{pmatrix}
  =
  \begin{pmatrix}
    \uvec{\i} & \uvec{\j}
  \end{pmatrix}
  \,
  \begin{pmatrix}
    1 & \cos\alpha\\
    0 & \sin\alpha
  \end{pmatrix}
\]

\subsubsection{Resumiendo el cambio de base}
Como dijimo, la matriz de cambio de base de $B$ a $B'$ está formada por las componentes
de los vectores de la expresión dispuestos en columnas. A continuación se representan, esta matriz de cambio de base y su inversa
\begin{subequations}
  \begin{align}
    \label{eq:R2-oblic_M}
    &\mmm{M}
      =
      \begin{pmatrix}
        1 & \cos\alpha\\
        0 & \sin\alpha
      \end{pmatrix}\\
    \label{eq:R2-oblic_Minv}
    &\mmm{M}^{-1}
      =
      \dfrac{1}{\sin\alpha}
      \begin{pmatrix}
        \sin\alpha & -\cos\alpha\\
        0 & 1        
      \end{pmatrix}
      =
      \begin{pmatrix}
        1 & -\cot\alpha\\
        0 & \phantom{-}\csc\alpha
      \end{pmatrix}
  \end{align}
\end{subequations}
Utilizaremos la matriz directa $M$ para transformar la base rectangular $B$ en la nueva
base oblícua $B'$. En cambio, para obtener la base rectangular a partir de la oblícua,
utilizaríamos la matriz inversa $\mmm{M}^{-1}$. Nótese que los vectores se representan en
filas
\begin{subequations}
\begin{align}
  &\begin{pmatrix}
    \xhat{e}_1 &  \xhat{e}_2
  \end{pmatrix}
  =
    \begin{pmatrix}
      \uvec{\i} & \uvec{\j}
    \end{pmatrix}
    \,
    \begin{pmatrix}
      1 & \cos\alpha\\
      0 & \sin\alpha
    \end{pmatrix}\\
  &\begin{pmatrix}
    \uvec{\i} & \uvec{\j}
  \end{pmatrix}
  =
    \begin{pmatrix}
      \xhat{e}_1 & \xhat{e}_2
    \end{pmatrix}
    \,
    \begin{pmatrix}
      1 & -\cot\alpha\\
      0 & \phantom{-}\csc\alpha
    \end{pmatrix}
    \,
\end{align}
\end{subequations}


\section{Componentes de los vectores}
Cuando un vector se expresa en una determinada base, se obtiene su representación en
función de unas determinadas componentes. Por ejemplo, en dos dimensiones y tomando una
base genérica $\set{\vvv{e}_1,\vvv{e}_2}$, escribiríamos
\[
  \vvv{v} = x_1\vvv{e}_1 + x_2\vvv{e}_2
\]
que, como sabemos, significa que el vector $\vvv{v}$ se obtiene sumando los productos de
cada cada componente por el vector correspondiente de la base.

Si la base estuviera formada por muchos vectores, o este número no estuviera definido,
entonces sería conveniente utilizar un sumatorio, como
\[
  \vvv{v} = \sum_{i=1}^n x_i\vvv{e}_i
\]

Cuando se opera con expresiones un tanto largas, se suele utilizar la notación de
Einstein, donde se prescinde del sumatorio y se tiene en cuenta que cuando en un
monomio se encuentra el mismo índice, por ejemplo $i$, se debe sumar para todos los
valores de ese índice
\[
  \vvv{v} = \sum_i x_i\vvv{e}_i = x_i \vvv{e}_i
\]

Otro aspecto importante es notar que las componentes de un vector dependen de la base
elegida. Si se cambia de base, las componentes varían. Por ejemplo, si utilizáramos
la base $\set{\vvv{x}'_1, \vvv{x}'_2}$, las nuevas componentes serían $x'_i$ y $x'_2$
\[ 
  \vvv{v} = x'_1\vvv{e}'_1 + x'_2\vvv{e}'_2
\]
Cambian las componentes, pero el vector tiene que ser el mismo, el vector no cambia, lo
hace su representación dependiendo de la base utilizada.

\section{Componentes contravariantes}
\label{ssect:R2-contravariantes}
Para introducir este tipo de componentes, vamos a poner un ejemplo concreto. Supongamos
que medimos un cierto vector y para ello, por sencillez, utilizamos la base estándar en coordenadas cartesianas rectangulares en metros $B=\set{\uvec{\i},\,\uvec{\j}}$, y
obtenemos que mide \qty{1}{\metre} en el eje $x$ y \qty{2}{\metre} en el eje $y$.
El vector se podrá representar como
\[
  \vvv{v} = 1\,\uvec{\i} + 2\,\uvec{\j}\,\unit{\metre}
\]
esto es, las componentes serían $(1, 2)\,\unit{\metre}$.

Supongamos también que deseamos cambiar la base elegida, a otra
$B' = \set{\uvec{\i}',\, \uvec{\j}'}$ que mide en centímetros. La nueva base
\emph{es más pequeña} que la antigua (centímetros en lugar de metros). Pero, para que el vector $\vvv{v}$ no cambie, los valores medidos en centímetros deben
\emph{aumentar} con respecto a las componentes en metros ($A'_i > A_i$)
\[
  \vvv{v} = 100\,\uvec{\i} + 200\,\uvec{\j}\,\unit{\centi\metre}
\]
las nuevas componentes serían $(100, 200)\,\unit{\centi\metre}$.
En otras palabras, la base ``se ha hecho 100 veces menor'', pero las componentes
``se han hecho 100 veces mayor''.

\subsubsection{Concepto}
Estas componentes han variado de forma contraria a como lo hizo la base elegida.
A las que se comportan de esta manera se les llama \emph{contravariantes} y, aunque no
es necesario en la mayor parte de las situaciones, se suele utilizar el índice que las
acompaña como \emph{superíndice} (que no representa un exponente).
Además, para facilitar la notación de Einstein, escribiremos las componentes
contravariantes como $x^i$ (y también los ejes correspondientes), en lugar de $x$ e $y$.
Así, un vector como el que hemos puesto de ejemplo, se podría representar en forma
genérica como
\[
  \vvv{v} = x^1\vvv{e}_1 + x^2\vvv{e}_2
\]
En notación de Einstein sería
\[
  \vvv{v} = x^i\vvv{e}_i
  \hspace{1em}\text{con } i=1,2\text{.}
\]
Los vectores de la base  $B'=\set{\vvv{e}_1, \vvv{e}_2}$, con el índice abajo, los
llamamos \emph{vectores covariantes} y decimos que forman una \emph{base covariante}.


\subsubsection{Representación gráfica de las componentes contravariantes}
La componentes contravariantes en el sistema de coordenadas oblicuas se pueden
representar gráficamente, trazando líneas paralelas a los ejes a partir del extremo del
segmento orientado $\vvv{v}$. Esta es la forma natural con la que sumamos vectores mediante la regla del paralelogramo, ver figura \ref{fig:R2-oblic_contrav}.
\begin{figure}[ht]
  % Escala
  \def\scl{1.00}
  % 
  %\pgfmathsetmacro{\EJESCARTLONG}{3.0}
  %\pgfmathsetmacro{\EJEOBLICLONG}{2.7}
  \pgfmathsetmacro{\EJESCARTLONG}{2.8}
  \pgfmathsetmacro{\EJEOBLICLONG}{2.5}
  % Eje x
  %\pgfmathsetmacro{\XLONG}{\EJESCARTLONG}
  % Eje y
  %\pgfmathsetmacro{\YLONG}{\EJESCARTLONG}
  % Eje xtilde
  \pgfmathsetmacro{\XTILDEANG}{0}
  \pgfmathsetmacro{\XTILDELONG}{\EJEOBLICLONG}    
  % Eje ytilde
  \pgfmathsetmacro{\YTILDEANG}{70}
  \pgfmathsetmacro{\YTILDELONG}{\EJESCARTLONG}
  % Vector P
  \pgfmathsetmacro{\PMOD}{2.3}
  \pgfmathsetmacro{\PANG}{40}
  % Fondo
  \pgfmathsetmacro{\HORZ}{0.25}
  \pgfmathsetmacro{\VERT}{0.25}
  % 
  \centering
 \begin{tikzpicture}[%
    scale=\scl,
    every node/.style={black,font=\small},
    %ejecart/.style={->, black!30},
    ejeoblic/.style={->, black},
    proyeccion/.style={black!20, line width=0.4},
    % longitud/.style={line width=.5pt,\colorV},
    texto/.style={black},
    %textoapagado/.style={\colorTorig},
    vector/.style={-{Latex[round]}, shorten >=1.2pt, line width=.8pt,\colorV},
    contravariante/.style={%
      -{Latex[round,width=5pt,length=8pt]},line width=1.2pt, orange!80!black},
    contravtxt/.style={orange!60!black},
    pcirculo/.style={fill=\colorPorig, draw=black},
    background/.style={%
      line width=\bgborderwidth,
      draw=\bgbordercolor,
      fill=\bgcolor,
    },
    ]
    % COORDENADAS
    % Origen
    \coordinate (O) at (0,0);
    % Eje xtilde
    \path[save path=\ejextilde] (O) -- ++(\XTILDEANG:\XTILDELONG cm)
    coordinate (xtildefin);
    \node[texto, below] at (xtildefin) {\footnotesize $x^1$};
    % Eje ytilde
    \path[save path=\ejeytilde, name path=O--ytilde] (O) -- ++(\YTILDEANG:\YTILDELONG cm)
    coordinate (ytildefin);
    % Vector
    \path[save path=\r] (O) -- coordinate[pos=0.6] (posvecttxt) (\PANG:\PMOD cm)
    coordinate (P);
    % Proyección de P en los ejes
    % Con eje x
    \path[save path=\proyx] (e2fin) |- (xfin) coordinate[pos=0.5] (e2x) (xfin);
    % Componentes contravariantes
    % Intersección de horizontal desde P hasta ytilde
    % -> Componente contravariante eje ytilde (CTV2) 
    \path[name path=P--y] (P) -| (yfin);
    \path[name intersections={of=P--y and O--ytilde,by=CTV2}];
    % Paralela (CTV2) -- (O) que pase por (P)
    % -> Componente contravariante eje xtilde (CTV1)
    \path (P) -- ($(P) - (CTV2) - (O)$) coordinate (CTV1);
    
    % DIBUJOS
    % Proyecciones de componentes
    % Contravariantes
    \draw[proyeccion] (P) -- (CTV1);
    %\node[below] at (CTV1) {\footnotesize $\hat{x}^1$};
    \draw[proyeccion] (P) -- (CTV2);
    %\node[left] at (CTV2) {\footnotesize $\hat{x}^2$};
    % xtilde
    \draw[ejeoblic, use path=\ejextilde];
    % ytilde
    \draw[ejeoblic, use path=\ejeytilde];
    \node[above right=0pt and -4pt, texto] at (ytildefin) {$x^2$};
    %\node[above right=0pt and -4pt, texto] at (ytildefin) {$\tilde{y}$};
    % Vector v
    \draw[vector, use path=\r];
    \node[above left=-3pt and -1pt] at (posvecttxt) {\footnotesize $\vvv{v}$};
    %\filldraw[fill=green, draw=black] (P) circle[radius=1.2pt];
    %\node[above right=0pt and -5pt] at (P) {\footnotesize $P$};

    % Marcar en color la componentes contravariantes
    \draw[contravariante]
    (O) -- node[contravtxt,pos=0.5, below] {$x^1\xhat{e}_1$} (CTV1);
    \draw[contravariante] (O) -- node[contravtxt,pos=0.5,above left=-2pt and -2pt]
    {$x^2\xhat{e}_2$} (CTV2);
    
    % Origen
    \fill[orange!80!black] (O) circle [radius=0.8pt];
    
    % Fondo amarillo
    \coordinate (SW) at ($(current bounding box.south west) + (-\HORZ cm,-\VERT cm)$);
    \coordinate (NE) at ($(current bounding box.north east) + (\HORZ cm,\VERT cm)$);
    \begin{scope}[on background layer]
      \draw[background] (SW) rectangle (NE);
    \end{scope}
  \end{tikzpicture}
  \caption{Componentes contravariantes $x^i$, en coordenadas oblicuas. El vector
    $\vvv{v}$ se representa como
    $\vvv{v}=x^1\,\xhat{e}_1+x^2\,\xhat{e}_2$. La componente
    contravariante en un cierto eje nos indica el desplazamiento que se ha producido en
    ese eje.}
  \label{fig:R2-oblic_contrav}
\end{figure}

Para centrar ideas, supongamos que el vector $\vvv{v}$ de la figura representa un
desplazamiento en $\symbb{R}^2$. La componente $x^1$ nos indica el desplazamiento
en el eje $x^1$, medido por el vector $\xhat{e}_1$ de la base, y $x^2$ nos lo da en el
eje $x^2$, medido por $\xhat{e}_2$.
La suma de las componentes contravariantes $x^i$ multiplicadas por sus
respectivos vectores covariantes de la base $\xhat{e}_i$, nos da el desplazamiento
indicado por el vector. Nótese que en la suma se utiliza la regla del paralelogramo.
Esto es característico de las componentes contravariantes que, en cierto modo
\emph{funcionan} como las componentes $x$ e $y$ en un sistema cartesiano ortogonal.


\section{Forma matricial de la base y de las componentes}
Recordemos que los vectores de una base se representan matricialmente situando los
vectores en fila
\[
  B =
  \begin{pmatrix}
    \vvv{e}_1 & \vvv{e}_2
  \end{pmatrix}
\]

En forma matricial, las componentes contravariantes del vector $\vvv{v}$, en una cierta
base, se representan en columna. Así, el vector anterior se puede expresar en la base
elegida como
\[
  \vvv{v} =
  \begin{pmatrix}
    x^1\\
    x^2
  \end{pmatrix}
\]
Es importante distinguir entre estas dos notaciones y darse cuenta de que no son
comparables. La primera, es la forma en la que se representa un conjunto de vectores que
forman una base, mientras que la segunda representa las componentes de un vector en una
cierta base. Lo único que tienen en común es que ambas son representaciones matriciales,
aunque de cosas diferentes.


\section{Transformación de componentes}
Un cambio de base, por ejemplo de $B$ a $B'$, se produce mediante una matriz de
transformación $\mmm{M}$
\begin{equation}
  B' = \mmm{M} B
\end{equation}
Entonces, las componentes de cualquier vector también se modifican.

Las componentes contravariantes de un vector se representan en columna y en la expresión
matemática de la transformación interviene la matriz inversa del cambio de base (porque
las componentes  contravariantes se transforman de forma inversa a como lo hacían los
vectores de la base). Así, las nuevas componentes se obtienen a partir de las antiguas
mediante $\mmm{M}^{-1}$, mientras que las antiguas se transforman a partir de las nuevas
mediante $\mmm{M}$
\begin{subequations}
  \begin{align}
    &\left[\tilde{x}^i\right]_{B'} = \mmm{M}^{-1} \left[x^i\right]_{B}\\
    &\left[x^i\right]_{B}\hspace{0.4em} = \mmm{M} \left[\tilde{x}^i\right]_{B'}
  \end{align}
\end{subequations}

En coordenadas oblicuas en $\symbb{R}^2$ utilizamos las matrices inversa
\eqref{eq:R2-oblic_Minv} y directa \eqref{eq:R2-oblic_M}, respectivamente
\begin{subequations}
  \begin{align}
    \begin{pmatrix}
      \tilde{x}^1\\
      \tilde{x}^2
    \end{pmatrix}
    &=
      \begin{pmatrix}
        1 & -\cot\alpha\\
        0 & \phantom{-}\csc\alpha
      \end{pmatrix}
      \,
    \begin{pmatrix}
      x^1\\
      x^2
    \end{pmatrix}\\
    \begin{pmatrix}
      x^1\\
      x^2
    \end{pmatrix}
    &=
      \begin{pmatrix}
        1 & \cos\alpha\\
        0 & \sin\alpha
      \end{pmatrix}
      \,
    \begin{pmatrix}
      \tilde{x}^1\\
      \tilde{x}^2
    \end{pmatrix}
  \end{align}
\end{subequations}


\subsubsection{Ejemplo}
Supongamos que tenemos un vector en coordenadas cartesianas rectangulares
\[
  \vvv{v}
  = x^1\,\xhat{u}_1 + x^2\,\xhat{u}_2 = x^i\,\xhat{u}_i
  = 2\,\uvec{\i} - 4\,\uvec{\j}
\]
donde $\xhat{u}_1 = \uvec{\i}$ y $\xhat{u}_2 = \uvec{\j}$.
Estamos interesados en obtener las componentes $\tilde{A}^i$ en un sistema de
coordenadas oblicuas, donde el ángulo que forman sus ejes vale
$\alpha = \pi/3\,\unit{\radian}$. Tendremos que aplicar la inversa de la matriz de
transformación de base de $B$ a $B'$
\begin{align*}
  \begin{pmatrix}
    \tilde{x}^1\\
    \tilde{x}^2
  \end{pmatrix}
  &=
  \begin{pmatrix}
    1 & -\cot{\pi/3}\\
    0 & \phantom{-}\csc{\pi/3}
  \end{pmatrix}
  \,
  \begin{pmatrix}
    x^1\\
    x^2
  \end{pmatrix}
    =
    \begin{pmatrix}
      1 & -1/\sqrt{3}\\
      0 & \phantom{-}2/\sqrt{3}
    \end{pmatrix}
    \,
    \begin{pmatrix}
      \phantom{-}2\\
      -4
    \end{pmatrix}
    =
    \begin{pmatrix}
      2+\dfrac{4}{\sqrt{3}}\\[3ex]
      -\dfrac{8}{\sqrt{3}}
    \end{pmatrix}
\end{align*}

Así
\[
  \vvv{v}
  = \tilde{x}^1\,\xhat{e}_1 + \tilde{x}^2\,\xhat{e}_2 = \tilde{x}^i\,\xhat{e}_i
  = \left(2+\dfrac{4}{\sqrt{3}}\right)\,\xhat{e}_1 - \dfrac{8}{\sqrt{3}}\,\xhat{e}_2
\]

Si nos dieran las componentes en coordenadas oblicuas
\[
  \left(\tilde{x}^1, \tilde{x}^2\right) = \left(2+4/\sqrt{3}, -8/\sqrt{3}\right)
\]
y quisiéramos obtener las coordenadas rectangulares, deberíamos obtener $(2,-4)$
\begin{align*}
  \begin{pmatrix}
    x^1\\
    x^2
  \end{pmatrix}
  &=
  \begin{pmatrix}
    1 & \cos{\pi/3}\\
    0 & \sin{\pi/3}
  \end{pmatrix}
  \,
  \begin{pmatrix}
    \tilde{x}^1\\
    \tilde{x}^2
  \end{pmatrix}
    =
  \begin{pmatrix}
    1 & 1/2\\
    0 & \sqrt{3}/2
  \end{pmatrix}
  \,
  \begin{pmatrix}
    2+4/\sqrt{3}\\
    -8/\sqrt{3}
  \end{pmatrix}\\
  &=
    \begin{pmatrix}
      2+4/\sqrt{3}-4/\sqrt{3}\\
      -\sqrt{3}/2\, \left(-8/\sqrt{3}\right)
    \end{pmatrix}
    =
    \begin{pmatrix}
      \phantom{-}2\\
      -4
    \end{pmatrix}  
\end{align*}


\section{Métrica}
Con los conocimientos que hemos adquirido, la forma más rápida para obtener el tensor
métrico en coordenadas oblicuas es utilizar el método seguido en la expresión
\eqref{eq:R2-cart-metrica}. Para ello debemos calcular el producto escalar de los
vectores de la base \eqref{eq:R2-oblic_base}.

Multiplicamos escalarmente estos vectores
\begin{subequations}
  \begin{align}
    \label{eq:R2-oblic-e1e1_prod_escalar}
    \xhat{e}_1\cdot\xhat{e}_1
    &=
      \uvec{\i}\cdot\uvec{\i} = 1\\
    \label{eq:R2-oblic-e1e2_prod_escalar}    
    \xhat{e}_1\cdot\xhat{e}_2
    &=
      \uvec{\i}\cdot (\cos\theta\,\uvec{i} + \sin\theta\,\uvec{\j}) = \cos\theta\\
    \label{eq:R2-oblic-e2e1_prod_escalar}
    \xhat{e}_2\cdot\xhat{e}_1
    &=
      \xhat{e}_1\cdot\xhat{e}_2 = \cos\theta\\
    \label{eq:R2-oblic-e2e2_prod_escalar}
    \xhat{e}_2\cdot\xhat{e}_2
    &=
      (\cos\theta\,\uvec{\i}+\sin\theta\,\uvec{j})
      \cdot (\cos\theta\,\uvec{\i}+\sin\theta\,\uvec{j})
      = (\cos^2\theta + \sin^2\theta) = 1
  \end{align}
\end{subequations}

El tensor métrico en coordenadas oblicuas resulta
\begin{equation}\label{eq:R2-oblic_metrica}
  \mmm{g}_{ij}
  = \begin{pmatrix}
    g_{11} & g_{12}\\
    g_{21} & g_{22}
    \end{pmatrix}
  = \begin{pmatrix}
      \xhat{e}_1\cdot\xhat{e}_1 & \xhat{e}_1\cdot\xhat{e}_2\\
      \xhat{e}_2\cdot\xhat{e}_1 & \xhat{e}_2\cdot\xhat{e}_2
  \end{pmatrix}
  = \begin{pmatrix}
    1 & \cos\alpha\\
    \cos\alpha & 1
  \end{pmatrix}
\end{equation}

Analizamos la métrica obtenida
\begin{itemize}
\item Los elementos diagonales valen la unidad. Esto significa que un cambio en
  cualquiera  de los ejes $\tilde{x}$ o $\tilde{y}$ o componentes $A^i$, representan ese
  mismo valor de desplazamiento en el espacio.
\item Los elementos no diagonales no son cero, lo que indica que un cambio cualquiera en
  las coordenadas implica un cambio en la otra, y viceversa (los ejes no son
  ortogonales).
  \begin{itemize}
  \item Un caso extremo posible se produciría cuando
    $\alpha = \pi/2\,\unit{\radian} = \qty{90}{\degree}$. En este caso
    $\cos\alpha = \cos\pi/2 = 0$ y el tensor métrico se representaría mediante la matriz
    unidad. En esta situación, los ejes serían ortogonales y no tendríamos un sistema de
    coordenadas oblículo, sino un sistema de coordenadas cartesianas rectangulares, como
    al inicio del capítulo.
  \item El otro caso extremo, que no produciría ningún sistema de referencia en
    $\symbb{R}^2$, se daría cuando $\alpha = 0$. En este caso $\cos\alpha = \cos 0 = 1$
    y los dos ejes $\tilde{x}$ y $\tilde{y}$, coincidirían. Además, los vectores
    unitarios respectivos, serían iguales $\xhat{e}_1 = \xhat{e}_2 = \uvec{\i}$ y ya no
    formarían una base del plano, ni tendríamos un sistema de coordenadas válido (el
    determinante de la métrica sería nulo, la matriz sería singular y no tendría
    inversa) y el sistema colapsaría. En todo caso tendríamos un subespacio vectorial de
    $\symbb{R}^2$, con una dimensión, definiendo el espacio de los reales $\symbb{R}$,
    pero no un sistema de coordenadas válido para el plano.
  \end{itemize}
\end{itemize}

Calculamos la métrica inversa
{\small
  \begin{align*}
    \mmm{g}^{ij}
    &=
      \dfrac{1}{\det(g_{ij})}
      \text{adj}\left(g_{ij}^\transp\right)
      = \dfrac{1}{\sin^2\alpha}
      \begin{pmatrix}
        1 & -\cos\alpha\\
        -\cos\alpha & 1
      \end{pmatrix}
      = \begin{pmatrix}
        \csc^2\alpha & -\csc\alpha\cot\alpha\\
        -\csc\alpha\cot\alpha & 1
      \end{pmatrix}
  \end{align*}
}
Resultando
{\small
\begin{equation}\label{eq:R2-oblic_metrica_inv}
  \mmm{g}^{ij}
  = \begin{pmatrix}
    g^{11} & g^{12}\\
    g^{21} & g^{22}
  \end{pmatrix}
  = \dfrac{1}{\sin^2\alpha}
  \begin{pmatrix}
    1 & -\cos\alpha\\
    -\cos\alpha & 1
  \end{pmatrix}
  = \begin{pmatrix}
    \csc^2\alpha & -\csc\alpha\cot\alpha\\
    -\csc\alpha\cot\alpha & 1
  \end{pmatrix}
\end{equation}
}

\section{Espacio tangente y espacio dual}
Si se quiere profundizar algo más en las componentes covariantes y contravariantes, aquí
se presentan algunas ideas útiles.

Hay una definición más formal de estas componentes, que procede del álgebra lineal y de
la geometría diferencial; de ambas, pero con propósitos y niveles de abstracción
distintos.
\begin{itemize}
\item Por parte del álgebra lineal, la distinción entre vectores y/o componentes
  covariantes y contravariantes nace de la relación entre un espacio vectorial $V$, y su
  espacio dual $V^*$:
\begin{itemize}
\item Vectores (contravariantes): Son los elementos del espacio original $V$. Sus
  elementos se transforman de manera inversa a los cambios de base, para que el vector
  físico permanezca igual.
\item Covectores o formas lineales (covariantes): Son elementos del espacio dual $V^*$
  asociado a $V$. Son funciones $\omega$ que actúan sobre los vectores y producen un
  escalar. Sus componentes se transforman de la misma forma que los cambios de base.
\end{itemize}
A este nivel, las componentes covariantes son simplemente los coeficientes de una
\emph{1-forma} o un funcional lineal ---las funciones que operan sobre un solo vector---.
Estos coeficientes \emph{viven} en un espacio dual $V^*$ asociado al espacio vectorial
$V$ original. En principio, esta es la relación entre los vectores y los covectores,
pero viven en espacios distintos. Ahora bien, si disponemos de una métrica (un producto
escalar definido en $V$, también lo tendremos en $V^*$) y ambos espacios se pueden
conectar y el álgebra lineal te permite \emph{bajar} o \emph{subir} índices convirtiendo
un vector en un covector y viceversa.

\item Por su parte, la geometría diferencial toma estos conceptos del álgebra lineal y
  los aplica a cada punto de una variedad ---una superficie curva, por ejemplo; una
  superficie esférica o tórica en un espacio $\symbb{R}^3$---.
  \begin{itemize}
  \item Espacio tangente: En cada punto $P$\, de una curva o superficie, existe una
    línea o un plano tangente al punto que es un espacio vectorial $V$, o
    $T_p\symbb{R}^2$. Aquí es donde viven los vectores y tiene sentido el álgebra
    lineal. En cada punto de la variedad hay definido un plano o espacio vectorial
    diferente. En principio, no podríamos operar o comparar vectores de diferentes
    espacios tangente (asociados a diferentes puntos de la variedad). Para ello, hay que
    definir los conceptos de \emph{transporte paralelo} y \emph{derivada covariante},
    para poder \emph{transportar} un vector desde su espacio en un punto a otro espacio
    asociado a otro punto, y con ello poder realizar las operaciones comunes, suma,
    resta, derivada, etc., en un mismo espacio.

    En nuestro caso, como $\symbb{R}^2$ es plano, pues no tiene más dimensiones, todo él
    es un mismo espacio vectorial, independientemente del punto de referencia y, por
    tanto, tienen sentido todas las operaciones de suma, resta, comparación, derivada,
    etc., sin necesidad del transporte paralelo y derivada covariante.
  \item Campos tensoriales: La geometría diferencial estudia cómo estas componentes
    (covariantes y contravariantes) cambian, no solo cuando cambias de base, sino cuando
    te mueves de un punto a otro en un sistema de coordenadas curvilíneo ---cuando se
    cambia de un punto a otro, los espacios tangentes cambian y, en principio, no se
    pueden relacionar fácilmente, pero la geometría diferencial nos permite comparar y
    operar con vectores de diferentes espacios tangente.

    En geometría diferencial, las componentes covariantes suelen estar asociadas a los
    gradientes de las funciones escalares (de ahí lo de las líneas o superficies de
    nivel, mientras que las contravariantes se asocian a velocidades o desplazamientos
    infinitesimales.
  \end{itemize}
\end{itemize}

Para terminar, el álgebra lineal proporciona la estructura algebraica, mientras que la geometría diferencial proporciona el marco geométrico y físico.

\section{Base contravariante}
Como dijimos en el apartado anterior, la métrica nos permite conectar la base covariante
$\set{\vvv{e}_1,\vvv{e}_2}$ con su dual, base contravariante, \emph{subiendo} índices.E
Utilizando la métrica inversa \eqref{eq:R2-oblic_metrica_inv}, y en notación de Einstein
\begin{equation}
  \vvv{e}^i = g^{ij} \vvv{e}_j
\end{equation}
Así
\begin{align}
  \vvv{e}^1
  &=
    g^{11}\vvv{e}_1 + g^{12}\vvv{e}_2
    = \dfrac{1}{\sin^2\alpha}
    \left(%
    \vvv{e}_1 - \cos\alpha\,\vvv{e}_2
    \right)\\
  \vvv{e}^2
  &=
    g^{21}\vvv{e}_1 + g^{22}\vvv{e}_2
    = \dfrac{1}{\sin^2\alpha}
    \left(%
    -\cos\alpha\,\vvv{e}_1 + \vvv{e}_2
    \right)
\end{align}

\subsubsection{Módulo de los elementos de la base dual}
Comprobamos el módulo de estos vectores y deduciremos que solo serían unitarios en un
sistema cartesiano rectangular $\alpha=\qty{90}{\degree}$
\begin{align}
  \left|\vvv{e}_1\right|
  &= \left(\dfrac{1}{\sin^2\alpha}\right)^2
    \sqrt{1^2 + (-\cos\alpha)^2}
    = \dfrac{\sqrt{1+\cos^2\alpha}}{\sin^4\alpha}\\
  \left|\vvv{e}_2\right|
  &= \left(\dfrac{1}{\sin^2\alpha}\right)^2
    \sqrt{(-\cos\alpha)^2 + 1}
    = \dfrac{\sqrt{1+\cos^2\alpha}}{\sin^4\alpha}
\end{align}

\section{Producto escalar de los vectores de las dos bases}
Empezamos demostrando que el producto de vectores de los distintos espacios con el mismo
índice da la unidad
\begin{equation}
  \vvv{e}^i \cdot \vvv{e}_i = 1
\end{equation}
\begin{align*}
  \vvv{e}^1\cdot\vvv{e}_1
  &=
    \dfrac{1}{\sin^2\alpha}
    \left(\vvv{e}_1-\cos\alpha\,\vvv{e}_2\right)\cdot\vvv{e}_1
    = \dfrac{1}{\sin^2\alpha}
    (\vvv{e}_1\cdot\vvv{e}_1 - \cos\alpha\,\vvv{e}_2\cdot\vvv{e}_1)\\
  &=
    \dfrac{1}{\sin^2\alpha}
    (g_{11}-\cos\alpha\,g_{21})
    = \dfrac{1}{\sin^2\alpha}(1-\cos^2\alpha)
    = \dfrac{1}{\sin^2\alpha}\sin^2\alpha
    = 1
\end{align*}
\begin{align*}
  \vvv{e}^2\cdot\vvv{e}_2
  &=
    \dfrac{1}{\sin^2\alpha}
    \left(-\cos\alpha\,\vvv{e}_1+\vvv{e}_2\right)\cdot\vvv{e}_2
    = \dfrac{1}{\sin^2\alpha}
    (-\cos\alpha\,\vvv{e}_1\cdot\vvv{e}_2+\vvv{e}_2\cdot\vvv{e}_2)
  = \\
  &= \dfrac{1}{\sin^2\alpha}
    (-\cos\alpha\,g_{12} + g_{22})
    = \dfrac{1}{\sin^2\alpha}(-\cos^2\alpha+1)
    = \dfrac{1}{\sin^2\alpha}\sin^2\alpha
    = 1
\end{align*}
Por otro lado, demostramos que los productos \emph{cruzados} son nulos
\begin{equation}
  \vvv{e}^i\cdot\vvv{e}_j = 0
  \hspace{1em}\text{con } i\neq j
\end{equation}
\begin{align*}
  \vvv{e}^1\cdot\vvv{e}_2
  &=
    \dfrac{1}{\sin^2\alpha}
    (\vvv{e}_1-\cos\alpha\,\vvv{e}_2)\cdot\vvv{e}_2
    = \dfrac{1}{\sin^2\alpha}
    (\vvv{e}_1\cdot\vvv{e}_2 - \cos\alpha\,\vvv{e}_2\cdot\vvv{e}_2)\\
  &=
    \dfrac{1}{\sin^2\alpha} (g_{12} - \cos\alpha\,g_{22})    
    = \dfrac{1}{\sin^2\alpha} (\cos\alpha - \cos\alpha)    
    = 0
\end{align*}
\begin{align*}
  \vvv{e}^2\cdot\vvv{e}_1
  &=
    \dfrac{1}{\sin^2\alpha}
    (-\cos\alpha\,\vvv{e}_1+\vvv{e}_2)\cdot\vvv{e}_1
    = \dfrac{1}{\sin^2\alpha}
    (-\cos\alpha\,\vvv{e}_1\cdot\vvv{e}_1 + \vvv{e}_2\cdot\vvv{e}_1)\\
  &=
    \dfrac{1}{\sin^2\alpha} (-\cos\alpha\,g_{11} + g_{21})
    = \dfrac{1}{\sin^2\alpha} (-\cos\alpha + \cos\alpha)    
    = 0
\end{align*}

Podemos resumir escribiendo
\begin{equation}
  \vvv{e}^i\cdot\vvv{e}_i
  = \vvv{e}_i\cdot\vvv{e}^i
  = \delta_{ij}
\end{equation}

Nótese que $e^2$ es ortogonal a $e_1$ y que $e^1$ lo es a $e_2$.
Esto significa que el eje $x^2$ es perpendicular a $x^1$ y que $x^1$ es perpendicular
a $x^2$ ---estamos hablando de los ejes de coordenadas, no de las componentes en esos
ejes---, ver figura \ref{fig:R2-oblic_ejes_ctv_cov}.

\begin{figure}[ht]
  % Escala
  \def\scl{1.00}
  % 
  \pgfmathsetmacro{\EJELONG}{3.3}
  % Eje x^1
  \pgfmathsetmacro{\XCTVANG}{0}
  % Eje x^2
  \pgfmathsetmacro{\YCTVANG}{60}
  %  Eje x_1
  \pgfmathsetmacro{\XCOVANG}{90}
  % Eje x_2
  \pgfmathsetmacro{\YCOVANG}{\YCTVANG-90}
  % Módulo unidad
  \pgfmathsetmacro{\MODULOUNO}{1.4}
  % Vector e_1
  \pgfmathsetmacro{\EXCOV}{\MODULOUNO}
  % Vector e_2
  \pgfmathsetmacro{\EYCOV}{\MODULOUNO}
  % Vector e^1
  \pgfmathsetmacro{\EXCTV}{\MODULOUNO/cos{\YCTVANG}}
  % Vector e^2
  \pgfmathsetmacro{\EYCTV}{\MODULOUNO/cos{\YCTVANG}}
  
%  % Lineas de nivel y
%  \pgfmathsetmacro{\YANGULOROT}{\YTILDEANG - 90}
%  \pgfmathsetmacro{\INICIOLONGITUD}{-0}
%  \pgfmathsetmacro{\YNUMLINEASMASUNA}{20}
%  \pgfmathsetmacro{\YNUMLINEAS}{\YNUMLINEASMASUNA - 1}
%  \pgfmathsetmacro{\YUNIT}{\XTILDELONG / \YNUMLINEAS}
%  \pgfmathsetmacro{\LONGITUDLINEAS}{\XTILDELONG - \INICIOLONGITUD - 0.25}
%  \pgfmathsetmacro{\THETAANG}{\YTILDEANG - \PANG}
%  \pgfmathsetmacro{\YPOSLINEACOV}{\PMOD*cos(\THETAANG)/(\YUNIT*cos(90-\YTILDEANG)}
  % Fondo
  \pgfmathsetmacro{\HORZ}{0.70}
  \pgfmathsetmacro{\VERT}{0.25}
  % 
  \centering
  \begin{tikzpicture}[%
    scale=\scl,
    every node/.style={black,font=\small},
    lineanivel/.style={black!20},
    ejectv/.style={->, black},
    ejecov/.style={->, red},
    proyeccion/.style={black!20, line width=0.4},
    textoctv/.style={black},
    textocov/.style={\colorTred},
    vectorcov/.style={-{Latex[round]}, shorten >=1.2pt, line width=1.5pt, black},
    vectorctv/.style={-{Latex[round]}, shorten >=1.2pt, line width=1.5pt, red!80!black},
    % covariante/.style={%
    % -{Latex[round,width=5pt,length=8pt]},line width=1.2pt, cyan!70!black},
    covariante/.style={line width=1.8pt, cyan!70!black},
    covtxt/.style={cyan!40!black},
    pcirculo/.style={fill=\colorPorig, draw=black},
    background/.style={%
      line width=\bgborderwidth,
      draw=\bgbordercolor,
      fill=\bgcolor,
    },
    ]
    % COORDENADAS
    % Origen
    \coordinate (O) at (0,0);
    % Eje x^1
    \path[save path=\ejexctv] (O) -- ++(\XCTVANG:\EJELONG cm) coordinate (xctvfin);
    % Eje x^2
    \path[save path=\ejeyctv] (O) -- ++(\YCTVANG:\EJELONG cm) coordinate (yctvfin);
    % Eje x_1
    \path[save path=\ejexcov] (O) -- ++(\XCOVANG:\EJELONG cm) coordinate (xcovfin);
    % Eje x_2
    \path[save path=\ejeycov] (O) -- ++(\YCOVANG:\EJELONG cm) coordinate (ycovfin);
    % Vectores
    % Vector e_1
    \path[save path=\excov] (O) -- ++(\XCTVANG:\EXCOV cm) coordinate (excovfin);
    % Vector e_2
    \path[save path=\eycov] (O) -- ++(\YCTVANG:\EYCOV cm) coordinate (eycovfin);
    % Vector e^1
    \path[save path=\exctv] (O) -- ++(\XCOVANG:\EXCTV cm) coordinate (exctvfin);
    % Vector e^2
    \path[save path=\eyctv] (O) -- ++(\YCOVANG:\EYCTV cm) coordinate (eyctvfin);
    %\path[save path=\ejeyctv, name path=O--ytilde] (O) -- ++(\YTILDEANG:\YTILDELONG cm)
%    coordinate (ytildefin);
%    % Vector
%    \path[save path=\r] (O) -- coordinate[pos=0.6] (posvecttxt) (\PANG:\PMOD cm)
%    coordinate (P);
%    % Componentes covariantes
%    % Proyección de P en los ejes
%    % Con eje x
%    \path[save path=\proyx] (e2fin) |- (xfin) coordinate[pos=0.5] (e2x) (xfin);
%    % Componente covariante con eje xtilde (COV1)
%    % \path (P) |- coordinate[pos=0.5] (COV1) (xtildefin);
%    % Perpendicular a O--ytilde que pasa por P
%    % -> Componente covariante con eje ytilde (COV2)
%    \path (P) -- ($(O)!(P)!(ytildefin)$) coordinate (COV2);
%    
    % DIBUJOS
%    % Covariantes
%    % \draw[proyeccion] (P) -- (COV1);
%    % \draw[proyeccion] (P) -- (COV2);
    
%    \foreach \num in {0,...,\YNUMLINEASMASUNA}
%    \draw[lineanivel,rotate=\YANGULOROT]
%    (\INICIOLONGITUD, \YUNIT * \num) -- ++(\LONGITUDLINEAS, 0);
%    
%    \foreach \num in {0,...,\YPOSLINEACOV}
%    \draw[lineanivel,rotate=\YANGULOROT,red!50]
%    (\INICIOLONGITUD, \YUNIT * \num) -- ++(\LONGITUDLINEAS, 0);      
%    
    % Ejes
    % x^1 = xctv
    \draw[ejectv, use path=\ejexctv];
    \node[textoctv, right] at (xctvfin) {\footnotesize $x^1$};
    % x^2 = yctv
    \draw[ejectv, use path=\ejeyctv];
    \node[textoctv, above right=0pt and -4pt] at (yctvfin) {$x^2$};
%    % \node[rotate=90, left=8pt] at (ytexto) {$r\sin\theta$};
    % x_1 = xcov
    \draw[ejecov, use path=\ejexcov];
    \node[textocov, above] at (xcovfin) {\footnotesize $x_1$};
    % x_2 = ycov
    \draw[ejecov, use path=\ejeycov];
    \node[textocov, below right=-3pt and -1pt] at (ycovfin) {\footnotesize $x_2$};
    % Vectores de la base
    % e_1 = excov
    \draw[vectorcov, use path=\excov];
    \node[textoctv, below=2pt] at (excovfin) {$\vvv{e}_1$};
    % e_2 = eycov
    \draw[vectorcov, use path=\eycov];
    \node[textoctv, above left=-2pt and 0pt] at (eycovfin) {$\vvv{e}_2$};
    % e^1 = exctv
    \draw[vectorctv, use path=\exctv];
    \node[textocov, left] at (exctvfin) {$\vvv{e}^1$};
    % e^2 = eyctv
    \draw[vectorctv, use path=\eyctv];
    \node[textocov, below left=0pt and 0pt] at (eyctvfin) {$\vvv{e}^2$};
    
    % Origen
    \fill (O) circle [radius=0.8pt];
    
    % Fondo amarillo
    \coordinate (SW) at ($(current bounding box.south west) + (-\HORZ cm,-\VERT cm)$);
    \coordinate (NE) at ($(current bounding box.north east) + (\HORZ cm,\VERT cm)$);
    \begin{scope}[on background layer]
      \draw[background] (SW) rectangle (NE);
    \end{scope}
  \end{tikzpicture}
  \caption{En negro, los ejes contravariantes $x^1$ y $x^2$, estos ejes son los que se
    utilizan para calcular las componentes contravariantes $x^1$ y $x^2$ de cualquier
    vector (el contexto ayuda a diferenciar ejes de componentes).
    En rojo, los ejes contravariantes. Nótese que cada $x_i$ es perpendicular a su
    homólogo $x^i$. Además se representan los vectores de las bases normal y dual.}
  \label{fig:R2-oblic_ejes_ctv_cov}
\end{figure}


\section{Componentes covariantes}
En la sección \ref{ssect:R2-contravariantes}, proyectamos un vector trazando líneas
paralelas a los ejes; así es como sumamos vectores gráficamente ---mediante la regla del
paralelogramo---. Los resultados de esas proyecciones paralelas eran sus componentes
contravariantes.

En cambio, hay otra forma de proyectar un vector, y es trazando líneas perpendiculares a
los ejes, figura \ref{fig:R2-oblic_cov}.
Estas \emph{componentes covariantes} las representaremos mediante subíndices $x_i$
\begin{equation}
  \vvv{v} = x_1\vvv{e}^1 + x_2\vvv{e}^2 = x_i \vvv{e}^i
\end{equation}
La base utilizada aquí $\vvv{e}^i$, no es la natural de coordenadas oblicuas,
\eqref{eq:R2-oblic_base}, pues contiene un superíndice (de nuevo, no es un exponente).

Mientras las componentes contravariantes te dicen cuánto \emph{avanzas} a lo largo
de cada eje, las covariantes miden la \emph{sombra} o proyección sobre ese eje.
En coordenadas cartesianas rectangulares ambas coinciden, pero en este sistema no.
En la figura \ref{fig:R2-oblic_contrav_cov} se pueden comparar visualmente ambas
componentes.
\begin{figure}[ht]
  \begin{minipage}{0.45\linewidth}  
    % Escala
    \def\scl{1.00}
    % 
    \pgfmathsetmacro{\EJESCARTLONG}{3.0}
    \pgfmathsetmacro{\EJEOBLICLONG}{2.7}
    % Eje xtilde
    \pgfmathsetmacro{\XTILDEANG}{0}
    \pgfmathsetmacro{\XTILDELONG}{\EJEOBLICLONG}    
    % Eje ytilde
    \pgfmathsetmacro{\YTILDEANG}{70}
    \pgfmathsetmacro{\YTILDELONG}{\EJESCARTLONG}
    % Vector P
    \pgfmathsetmacro{\PMOD}{2.3}
    \pgfmathsetmacro{\PANG}{40}
    % Fondo
    \pgfmathsetmacro{\HORZ}{0.90}
    \pgfmathsetmacro{\VERT}{0.25}
    % 
    \centering
    \begin{tikzpicture}[%
      scale=\scl,
      every node/.style={black,font=\small},
      ejeoblic/.style={->, black},
      proyeccion/.style={black!20, line width=0.4},
      texto/.style={black},
      vector/.style={-{Latex[round]}, shorten >=1.2pt, line width=.8pt,\colorV},
      covariante/.style={%
        -{Latex[round,width=5pt,length=8pt]},line width=1.2pt, cyan!70!black},
      covtxt/.style={cyan!40!black},
      pcirculo/.style={fill=\colorPorig, draw=black},
      background/.style={%
        line width=\bgborderwidth,
        draw=\bgbordercolor,
        fill=\bgcolor,
      },
      ]
      % COORDENADAS
      % Origen
      \coordinate (O) at (0,0);
      % Eje xtilde
      \path[save path=\ejextilde] (O) -- ++(\XTILDEANG:\XTILDELONG cm)
      coordinate (xtildefin);
      % Eje ytilde
      \path[save path=\ejeytilde,name path=O--ytilde](O) -- ++(\YTILDEANG:\YTILDELONG cm)
      coordinate (ytildefin);
      % Vector
      \path[save path=\r] (O) -- coordinate[pos=0.6] (posvecttxt) (\PANG:\PMOD cm)
      coordinate (P);
      % Componentes covariantes
      % Proyección de P en los ejes
      % Con eje x
      \path[save path=\proyx] (e2fin) |- (xfin) coordinate[pos=0.5] (e2x) (xfin);
      % Componente covariante con eje xtilde (COV1)
      \path (P) |- coordinate[pos=0.5] (COV1) (xtildefin);
      % Perpendicular a O--ytilde que pasa por P
      % -> Componente covariante con eje ytilde (COV2)
      \path (P) -- ($(O)!(P)!(ytildefin)$) coordinate (COV2);
      
      % DIBUJOS
      % Covariantes
      \draw[proyeccion] (P) -- (COV1);
      \draw[proyeccion] (P) -- (COV2);
      
      % Ejes
      % xtilde
      \draw[ejeoblic, use path=\ejextilde];
      % \node[texto, right] at (xtildefin) {\footnotesize $\tilde{x}^1$};
      \node[texto, right] at (xtildefin) {\footnotesize $\tilde{x}$};
      % ytilde
      \draw[ejeoblic, use path=\ejeytilde];
      % \node[above right=0pt and -4pt, texto] at (ytildefin) {$\tilde{x}^2$};
      \node[texto, above right=0pt and -4pt] at (ytildefin) {$\tilde{y}$};
      % \node[rotate=90, left=8pt] at (ytexto) {$r\sin\theta$};
      % Vector P
      \draw[vector, use path=\r];
      \node[above left=-3pt and -1pt] at (posvecttxt) {\footnotesize $\vvv{v}$};
      % \filldraw[fill=green, draw=black] (P) circle[radius=1.2pt];
      % \node[above right=0pt and -5pt] at (P) {\footnotesize $P$};
      
      % Marcar en azul la componentes covariantes
      \draw[covariante] (O) -- node[covtxt,pos=0.5, below] {$x_1\vvv{e}^1$} (COV1);
      \draw[covariante] (O) -- node[covtxt,pos=0.5,above left=-2pt and -2pt]
      {$x_2\vvv{e}^2$} (COV2);
      
      % Origen
      \fill[cyan!80!black] (O) circle [radius=0.8pt];
      
      % Fondo amarillo
      \coordinate (SW) at ($(current bounding box.south west) + (-\HORZ cm,-\VERT cm)$);
      \coordinate (NE) at ($(current bounding box.north east) + (\HORZ cm,\VERT cm)$);
      \begin{scope}[on background layer]
        \draw[background] (SW) rectangle (NE);
      \end{scope}
    \end{tikzpicture}
    \caption{Componentes covariantes $x_i$ en coordenadas oblicuas. El vector $\vvv{v}$
      se representa como $\vvv{v}=x_1\,\vvv{e}^1 + x_2\,\vvv{e}^2$. Las componentes
      contravariantes nos indican el desplazamiento en el espacio.}
    \label{fig:R2-oblic_cov}
  \end{minipage}
  \hspace{1em}
  \begin{minipage}{0.45\linewidth}
    % Escala
    \def\scl{1.00}
    % 
    \pgfmathsetmacro{\EJESCARTLONG}{3.0}
    \pgfmathsetmacro{\EJEOBLICLONG}{2.7}
    % Eje x
    \pgfmathsetmacro{\XLONG}{\EJESCARTLONG}
    % Eje y
    \pgfmathsetmacro{\YLONG}{\EJESCARTLONG}
    % Eje xtilde
    \pgfmathsetmacro{\XTILDEANG}{0}
    \pgfmathsetmacro{\XTILDELONG}{\EJEOBLICLONG}    
    % Eje ytilde
    \pgfmathsetmacro{\YTILDEANG}{70}
    \pgfmathsetmacro{\YTILDELONG}{\EJESCARTLONG}
    % Vector P
    \pgfmathsetmacro{\PMOD}{2.3}
    \pgfmathsetmacro{\PANG}{40}
    % Fondo
    \pgfmathsetmacro{\HORZ}{0.90}
    \pgfmathsetmacro{\VERT}{0.25}
    % 
    \centering
    \begin{tikzpicture}[%
      scale=\scl,
      every node/.style={black,font=\small},
      ejecart/.style={->, black!30},
      ejeoblic/.style={->, black},
      proyeccion/.style={black!20, line width=0.4},
      % longitud/.style={line width=.5pt,\colorV},
      texto/.style={black},
      textoapagado/.style={\colorTorig},
      vector/.style={-{Latex[round]}, shorten >=1.2pt, line width=.8pt,\colorV},
      % vtilde/.style={%
      % -{Latex[round,width=4pt,length=6pt]}, line width=1pt, \colorVred},
      pcirculo/.style={fill=\colorPorig, draw=black},
      angulo/.style={fill=green!15, draw=\colorGirodos},
      background/.style={%
        line width=\bgborderwidth,
        draw=\bgbordercolor,
        fill=\bgcolor,
      },
      ]
      % COORDENADAS
      % Origen
      \coordinate (O) at (0,0);
      % Eje x
      \path[save path=\ejex] (O) -- (\XLONG cm, 0) coordinate (xfin);
      % Eje y
      \path[save path=\ejey] (O) -- (0, \YLONG cm) coordinate (yfin);
      % Eje xtilde
      \path[save path=\ejextilde] (O) -- ++(\XTILDEANG:\XTILDELONG cm)
      coordinate (xtildefin);
      \node[texto, below] at (xtildefin) {\footnotesize $\tilde{x}$};
      % Eje ytilde
      \path[save path=\ejeytilde,name path=O--ytilde](O) -- ++(\YTILDEANG:\YTILDELONG cm)
      coordinate (ytildefin);
      % Vector
      \path[save path=\r] (O) -- coordinate[pos=0.5] (pmid) (\PANG:\PMOD cm)
      coordinate (P);
      % Proyección de P en los ejes
      % Con eje x
      \path[save path=\proyx] (e2fin) |- (xfin) coordinate[pos=0.5] (e2x) (xfin);
      % Posición de texto en eje x
      % \path (O) -- coordinate[pos=0.5] (extxt) (e2x);
      % Con eje y
      \path[save path=\proyy] (e2fin) -| coordinate[pos=0.5] (e2y) (yfin);
      % Posición de texto en eje y
      % \path (O) -- coordinate[pos=0.5] (eytxt) (e2y);
      % Componentes covariantes y contravariantes
      % Intersección de horizontal desde P hasta ytilde
      % -> Componente contravariante eje ytilde (CTV2) 
      \path[name path=P--y] (P) -| (yfin);
      \path[name intersections={of=P--y and O--ytilde,by=CTV2}];
      % Paralela (CTV2) -- (O) que pase por (P)
      % -> Componente contravariante eje xtilde (CTV1)
      \path (P) -- ($(P) - (CTV2) - (O)$) coordinate (CTV1);
      % Componente covariante con eje xtilde (COV1)
      \path (P) |- coordinate[pos=0.5] (COV1) (xtildefin);
      % Perpendicular a O--ytilde que pasa por P
      % -> Componente covariante con eje ytilde (COV2)
      \path (P) -- ($(O)!(P)!(ytildefin)$) coordinate (COV2);
      % Proyecciones de las componentes
      
      % DIBUJOS
      % Ángulo theta
      \path (xfin) -- (O) -- (ytildefin) pic
      [angulo, angle radius=15pt]
      {angle = xfin--O--ytildefin};
      % \node[above right=-2pt and 4pt] at (O) {\scriptsize $\alpha$};
      % 
      % Proyecciones de componentes
      % Contravariantes
      \draw[proyeccion] (P) -- (CTV1);
      \node[orange!70!black, below=-1pt] at (CTV1) {\footnotesize $x^1$};
      \draw[proyeccion] (P) -- (CTV2);
      \node[orange!70!black, left=-1.5pt] at (CTV2) {\footnotesize $x^2$};
      % Covariantes
      \draw[proyeccion] (P) -- (COV1);
      \node[cyan!50!black, below=1pt] at (COV1) {\footnotesize $x_1$};
      \draw[proyeccion] (P) -- (COV2);
      \node[cyan!50!black, above left=-2pt and -2pt] at (COV2) {\footnotesize $x_2$};
      
      % Ejes
      % x
      \draw[ejecart, use path=\ejex];
      \node[textoapagado, below=1.2pt] at (xfin) {\small $x$};
      % y
      \draw[ejecart, use path=\ejey];
      \node[above, textoapagado] at (yfin) {$y$};
      % xtilde
      \draw[ejeoblic, use path=\ejextilde];
      % \node[below, texto] at (xtildefin) {$\tilde{x}^1$};
      % ytilde
      \draw[ejeoblic, use path=\ejeytilde];
      \node[above right=0pt and -4pt, texto] at (ytildefin) {$\tilde{y}$};
      % \node[rotate=90, left=8pt] at (ytexto) {$r\sin\theta$};
      % Vector P
      \draw[vector, use path=\r];
      \node[above left=-3pt and -1pt] at (pmid) {\footnotesize $\vvv{r}$};
      \filldraw[fill=green, draw=black] (P) circle[radius=1.2pt];
      \node[above right=0pt and -5pt] at (P) {\footnotesize $P$};
      
      % Origen
      \filldraw (O) circle [radius=.2pt];
      
      % Fondo amarillo
      \coordinate (SW) at ($(current bounding box.south west) + (-\HORZ cm,-\VERT cm)$);
      \coordinate (NE) at ($(current bounding box.north east) + (\HORZ cm,\VERT cm)$);
      \begin{scope}[on background layer]
        \draw[background] (SW) rectangle (NE);
      \end{scope}
    \end{tikzpicture}
    \caption{Componentes contravariantes, en naranja y covariantes, en azul. El ángulo
      entre los ejes está marcado en verde. El sistema de coordenadas rectangular se
      muestra, algo apagado, como referencia.}
    \label{fig:R2-oblic_contrav_cov}
  \end{minipage}
\end{figure}
Piensa en la base $\vvv{e}^i$ como una serie de líneas o superficies de nivel perpendiculares a cada uno de los ejes, que nos indican el valor de un cierto escalar o \emph{altura}. La componente covariante en ese eje nos indica de forma visual el
\emph{número de líneas o superficies de nivel} que el vector \emph{atraviesa}.
Ver figuras \ref{fig:R2-oblic_cov1} y \ref{fig:R2-oblic_cov2} (comparar con figura
\ref{fig:R2-oblic_ejes_ctv_cov}).
\begin{figure}[ht]
  \begin{minipage}{0.45\linewidth}
    % Escala
    \def\scl{1.00}
    % 
    % \pgfmathsetmacro{\EJESCARTLONG}{3.0}
    % \pgfmathsetmacro{\EJEOBLICLONG}{2.7}
     \pgfmathsetmacro{\EJESCARTLONG}{3.3}
    \pgfmathsetmacro{\EJEOBLICLONG}{3.0}
    %Eje xtilde
    \pgfmathsetmacro{\XTILDEANG}{0}
    \pgfmathsetmacro{\XTILDELONG}{\EJEOBLICLONG}
    % Eje ytilde
    \pgfmathsetmacro{\YTILDEANG}{70}
    \pgfmathsetmacro{\YTILDELONG}{\EJESCARTLONG}
    % Vector P
    \pgfmathsetmacro{\PMOD}{2.2}
    \pgfmathsetmacro{\PANG}{40}
    \pgfmathsetmacro{\XVECTOR}{\PMOD * cos(\PANG)}
    % Líneas de nivel x
    \pgfmathsetmacro{\INICIOALTURA}{-0}
    \pgfmathsetmacro{\XNUMLINEASMASUNA}{20}
    \pgfmathsetmacro{\XNUMLINEAS}{\XNUMLINEASMASUNA - 1}
    \pgfmathsetmacro{\XUNIT}{\XTILDELONG / \XNUMLINEASMASUNA}
    \pgfmathsetmacro{\ALTURALINEAS}{\YTILDELONG*cos(90-\YTILDEANG)-\INICIOALTURA - 0.1}
    \pgfmathsetmacro{\XPOSLINEACOV}{\XVECTOR / \XUNIT}
    % Fondo
    \pgfmathsetmacro{\HORZ}{0.93}
    \pgfmathsetmacro{\VERT}{0.44}
    % 
    \centering
    \begin{tikzpicture}[%
      scale=\scl,
      every node/.style={black,font=\small},
      lineanivel/.style={black!20},
      ejeoblic/.style={->, black},
      proyeccion/.style={black!20, line width=0.4},
      texto/.style={black},
      vector/.style={-{Latex[round]}, shorten >=1.2pt, line width=.8pt, black!40},
     % covariante/.style={%
     % -{Latex[round,width=5pt,length=8pt]},line width=1.2pt, cyan!70!black},
     covariante/.style={line width=1.8pt, cyan!70!black},
      covtxt/.style={cyan!40!black},
      pcirculo/.style={fill=\colorPorig, draw=black},
      background/.style={%
        line width=\bgborderwidth,
        draw=\bgbordercolor,
        fill=\bgcolor,
      },
      ]
      % COORDENADAS
      % Origen
      \coordinate (O) at (0,0);
      % Eje xtilde
      \path[save path=\ejextilde] (O) -- ++(\XTILDEANG:\XTILDELONG cm)
      coordinate (xtildefin);
      % Eje ytilde
      \path[save path=\ejeytilde, name path=O--ytilde]
      (O) -- ++(\YTILDEANG:\YTILDELONG cm)
      coordinate (ytildefin);
      % Vector
      \path[save path=\r] (O) -- coordinate[pos=0.6] (posvecttxt) (\PANG:\PMOD cm)
      coordinate (P);
      % Componentes covariantes
      % Proyección de P en los ejes
      % Con eje x
      \path[save path=\proyx] (e2fin) |- (xfin) coordinate[pos=0.5] (e2x) (xfin);
      % Componente covariante con eje xtilde (COV1)
      \path (P) |- coordinate[pos=0.5] (COV1) (xtildefin);
      % Perpendicular a O--ytilde que pasa por P
      % -> Componente covariante con eje ytilde (COV2)
      \path (P) -- ($(O)!(P)!(ytildefin)$) coordinate (COV2);
    
      % DIBUJOS
      % Covariantes
      %\draw[proyeccion] (P) -- (COV1);
      %\draw[proyeccion] (P) -- (COV2);

      \foreach \num in {0,...,\XNUMLINEAS}
      \draw[lineanivel] (\XUNIT * \num, \INICIOALTURA) -- ++(0, \ALTURALINEAS);

      \foreach \num in {0,...,\XPOSLINEACOV}
      \draw[lineanivel,red!50] (\XUNIT * \num, \INICIOALTURA) -- ++(0, \ALTURALINEAS);
      
      % Ejes
      % xtilde
      \draw[ejeoblic, use path=\ejextilde];
      %\node[texto, right] at (xtildefin) {\footnotesize $\tilde{x}^1$};
      \node[texto, right] at (xtildefin) {\footnotesize $\tilde{x}$};
      % ytilde
      \draw[ejeoblic, use path=\ejeytilde];
      %\node[above right=0pt and -4pt, texto] at (ytildefin) {$\tilde{x}^2$};
      \node[texto, above right=0pt and -4pt] at (ytildefin) {$\tilde{y}$};
      % \node[rotate=90, left=8pt] at (ytexto) {$r\sin\theta$};
      % Vector P
      \draw[vector, use path=\r];
      \node[above left=-3pt and -1pt,gray] at (posvecttxt) {\footnotesize $\vvv{v}$};
      %\filldraw[fill=green, draw=black] (P) circle[radius=1.2pt];
      %\node[above right=0pt and -5pt] at (P) {\footnotesize $P$};

      % Marcar en azul la componentes covariantes
      %\draw[covariante]
      %(O) -- node[covtxt,pos=0.5, below]
      %{$A_1\vvv{e}^1$} (COV1);
      \draw[covariante] (O) --node[covtxt,pos=0.5,below] {$x_1$} (COV1);
      %\draw[covariante] (O) -- node[covtxt,pos=0.5,above left=-2pt and -2pt]
      %{$\hat{x}_2\vvv{e}^2$} (COV2);
      
      % Origen
      \fill[cyan!80!black] (O) circle [radius=0.8pt];
    
      % Fondo amarillo
      \coordinate (SW) at ($(current bounding box.south west) + (-\HORZ cm,-\VERT cm)$);
      \coordinate (NE) at ($(current bounding box.north east) + (\HORZ cm,\VERT cm)$);
      \begin{scope}[on background layer]
        \draw[background] (SW) rectangle (NE);
      \end{scope}
    \end{tikzpicture}
    \caption{En azul, la componente covariante $x_1$ de un vector $\vvv{v}$ en
      coordenadas oblicuas, interpretada geométricamente como la proyección del vector o
      su \emph{sombra} en el eje $\tilde{x}$. En rojo aparecen las líneas de nivel
      perpendiculares al eje que son \emph{atravesadas} por $\vvv{v}$.
      En gris, las líneas que no son \emph{atravesadas}.}
    \label{fig:R2-oblic_cov1}
  \end{minipage}
  \hspace{1em}
  \begin{minipage}{0.45\linewidth}
    % Escala
    \def\scl{1.00}
    % 
    %\pgfmathsetmacro{\EJESCARTLONG}{3.0}
    %\pgfmathsetmacro{\EJEOBLICLONG}{2.7}
    \pgfmathsetmacro{\EJESCARTLONG}{3.3}
    \pgfmathsetmacro{\EJEOBLICLONG}{3.0}
    % Eje xtilde
    \pgfmathsetmacro{\XTILDEANG}{0}
    \pgfmathsetmacro{\XTILDELONG}{\EJEOBLICLONG}
    % Eje ytilde
    \pgfmathsetmacro{\YTILDEANG}{70}
    \pgfmathsetmacro{\YTILDELONG}{\EJESCARTLONG}
    % Vector P
    \pgfmathsetmacro{\PMOD}{2.22}
    \pgfmathsetmacro{\PANG}{40}
    \pgfmathsetmacro{\YVECTOR}{\PMOD * sin(\PANG)}
    % Lineas de nivel y
    \pgfmathsetmacro{\YANGULOROT}{\YTILDEANG - 90}
    \pgfmathsetmacro{\INICIOLONGITUD}{-0}
    \pgfmathsetmacro{\YNUMLINEASMASUNA}{20}
    \pgfmathsetmacro{\YNUMLINEAS}{\YNUMLINEASMASUNA - 1}
    \pgfmathsetmacro{\YUNIT}{\XTILDELONG / \YNUMLINEAS}
    \pgfmathsetmacro{\LONGITUDLINEAS}{\XTILDELONG - \INICIOLONGITUD - 0.25}
    \pgfmathsetmacro{\THETAANG}{\YTILDEANG - \PANG}
    \pgfmathsetmacro{\YPOSLINEACOV}{\PMOD*cos(\THETAANG)/(\YUNIT*cos(90-\YTILDEANG)}
    % Fondo
    \pgfmathsetmacro{\HORZ}{0.70}
    \pgfmathsetmacro{\VERT}{0.25}
    % 
    \centering
    \begin{tikzpicture}[%
      scale=\scl,
      every node/.style={black,font=\small},
      lineanivel/.style={black!20},
      ejeoblic/.style={->, black},
      proyeccion/.style={black!20, line width=0.4},
      texto/.style={black},
      vector/.style={-{Latex[round]}, shorten >=1.2pt, line width=.8pt, black!40},
      %covariante/.style={%
      %  -{Latex[round,width=5pt,length=8pt]},line width=1.2pt, cyan!70!black},
      covariante/.style={line width=1.8pt, cyan!70!black},
      covtxt/.style={cyan!40!black},
      pcirculo/.style={fill=\colorPorig, draw=black},
      background/.style={%
        line width=\bgborderwidth,
        draw=\bgbordercolor,
        fill=\bgcolor,
      },
      ]
      % COORDENADAS
      % Origen
      \coordinate (O) at (0,0);
      % Eje xtilde
      \path[save path=\ejextilde] (O) -- ++(\XTILDEANG:\XTILDELONG cm)
      coordinate (xtildefin);
      
      % Eje ytilde
      \path[save path=\ejeytilde, name path=O--ytilde]
      (O) -- ++(\YTILDEANG:\YTILDELONG cm)
      coordinate (ytildefin);
      % Vector
      \path[save path=\r] (O) -- coordinate[pos=0.6] (posvecttxt) (\PANG:\PMOD cm)
      coordinate (P);
      % Componentes covariantes
      % Proyección de P en los ejes
      % Con eje x
      \path[save path=\proyx] (e2fin) |- (xfin) coordinate[pos=0.5] (e2x) (xfin);
      % Componente covariante con eje xtilde (COV1)
      %\path (P) |- coordinate[pos=0.5] (COV1) (xtildefin);
      % Perpendicular a O--ytilde que pasa por P
      % -> Componente covariante con eje ytilde (COV2)
      \path (P) -- ($(O)!(P)!(ytildefin)$) coordinate (COV2);
    
      % DIBUJOS
      % Covariantes
      %\draw[proyeccion] (P) -- (COV1);
      %\draw[proyeccion] (P) -- (COV2);
      
      \foreach \num in {0,...,\YNUMLINEASMASUNA}
      \draw[lineanivel,rotate=\YANGULOROT]
      (\INICIOLONGITUD, \YUNIT * \num) -- ++(\LONGITUDLINEAS, 0);

      \foreach \num in {0,...,\YPOSLINEACOV}
      \draw[lineanivel,rotate=\YANGULOROT,red!50]
      (\INICIOLONGITUD, \YUNIT * \num) -- ++(\LONGITUDLINEAS, 0);      
      
      % Ejes
      % xtilde
      \draw[ejeoblic, use path=\ejextilde];
      %\node[texto, right] at (xtildefin) {\footnotesize $\tilde{x}^1$};
      \node[texto, right] at (xtildefin) {\footnotesize $\tilde{x}$};
      % ytilde
      \draw[ejeoblic, use path=\ejeytilde];
      %\node[above right=0pt and -4pt, texto] at (ytildefin) {$\tilde{x}^2$};
      \node[texto, above right=0pt and -4pt] at (ytildefin) {$\tilde{y}$};
      % \node[rotate=90, left=8pt] at (ytexto) {$r\sin\theta$};
      % Vector P
      \draw[vector, use path=\r];
      \node[above left=-3pt and -1pt,gray] at (posvecttxt) {\footnotesize $\vvv{v}$};
      %\filldraw[fill=green, draw=black] (P) circle[radius=1.2pt];
      %\node[above right=0pt and -5pt] at (P) {\footnotesize $P$};

      % Marcar en azul la componentes covariantes
      %\draw[covariante]
      %(O) -- node[covtxt,pos=0.5, below]
      %{$\hat{A}_1\xhat{e}^1$} (COV1);
      \draw[covariante] (O)--node[covtxt,pos=0.5,above left=-2pt and -2pt]{$x_2$} (COV2);
      
      % Origen
      \fill[cyan!80!black] (O) circle [radius=0.8pt];
    
      % Fondo amarillo
      \coordinate (SW) at ($(current bounding box.south west) + (-\HORZ cm,-\VERT cm)$);
      \coordinate (NE) at ($(current bounding box.north east) + (\HORZ cm,\VERT cm)$);
      \begin{scope}[on background layer]
        \draw[background] (SW) rectangle (NE);
      \end{scope}
    \end{tikzpicture}
    \caption{En azul, la componente covariante $x_2$ de un vector $\vvv{v}$ en
      coordenadas oblicuas, interpretada geométricamente como la proyección del vector o
      su \emph{sombra} en el eje $\tilde{y}$. En rojo aparecen las líneas de nivel
      perpendiculares al eje que son \emph{atravesadas} por $\vvv{v}$.
      En gris, las líneas que no son \emph{atravesadas}.}
    \label{fig:R2-oblic_cov2}
  \end{minipage}
\end{figure}

Una forma de asimilar este tipo de componentes es suponer que un vector $\vvv{v}$
representa el gradiente de un cierto campo escalar $\vvv{\nabla}\phi$.
Este gradiente se puede representar gráficamente como una serie de líneas de contorno o
superficies de nivel, que indica cuánto cambia un determinado escalar cuando nos
desplazamos en el espacio (en nuestro caso concreto en $\symbb{R}^2$).

Si aumentáramaos la escala de medida (utilizando una base \emph{mayor}, por ejemplo, en
el caso de distancias, usar metros en lugar de centímetros), las superficies de nivel
(gradiente) se comprimirían y, aunque el vector se vería más pequeño proporcionalmente
(componentes contravariantes), la densidad de líneas de nivel aumentaría, también
proporcionalmente.

En otras palabras, cuando un vector se visualiza mediante un cierto tipo de
recorrido en el espacio, como desplazamiento, velocidad, aceleración, etc.
---mediante sus componentes contravariantes---, al aumentar la escala de medida, las
líneas de nivel se comprimen, aumenta su densidad y el cambio en la ``altura'' debido al
gradiente es mayor, las componentes covariantes aumentan en consecuencia. Estas
componentes no varían inversamente con la escala de medida, sino que lo hace en el mismo
sentido. Estas son las \emph{componentes covariantes}.


\section{Producto escalar}
Ahora que se han introducido las componentes covariantes y contravariantes, y el tensor
métrico, estamos en disposición de ver el producto escalar de dos vectores
contravariantes.

Definiremos el producto escalar de dos vectores $\vvv{u}$ y $\vvv{v}$ matricialmente en
$\symbb{R}^2$ como
\begin{equation}
  \vvv{u}\cdot\vvv{v}
  = \vvv{u}^\transp  \mmm{g}_{ij} \vvv{v}
  = \begin{pmatrix}
    u^1 & u^2
  \end{pmatrix}
  \,\begin{pmatrix}
    g_{11} & g_{12}\\
    g_{21} & g_{22}
  \end{pmatrix}
  \,
  \begin{pmatrix}
    v^1\\
    v^2
  \end{pmatrix}
\end{equation}

Desarrollando la expresión
{\small
  \begin{align*}
    \vvv{u}\cdot\vvv{v}
    &=
      \begin{pmatrix}
        u^1 g_{11}+ u^2g_{21} & u^1g_{12}+ u^2g_{22}
      \end{pmatrix}
      \,
      \begin{pmatrix}
        v^1\\
        v^2
      \end{pmatrix}\\
    &=
      \left(u^1\,g_{11}+ u^2\,g_{21}\right) v^1+\left(u^1\,g_{12}+u^2\,g_{22}\right) v^2\\
    &=
      u^1g_{11}\,v^1 + u^2g_{21}\,v^1 + u^1g_{12}\,v^2 + u^2g_{22}\,v^2\\
    &=
      u^1\left(g_{11}\,v^1+g_{12}\,u^2\right) + v^2\left(g_{21}\,v^1+g_{22}\,v^2\right)\\
    &=
      u^1\left(\sum_{j=1}^2g_{1j}\,v^j\right) + u^2\left(\sum_{j=1}^2g_{2j}\,v^j\right)\\
    &=
      \sum_{i=1}^2 u^i \left(\sum_{j=1}^2g_{ij}\,v^j\right)\\
    &=
      \sum_{i=1}^2\sum_{j=1}^2g_{ij}\,u^i\,v^j
  \end{align*}
}

En notación de Einstein
\begin{equation}
  \vvv{u}\cdot\vvv{v}
  = g_{ij}\,u^i\,v^j
\end{equation}

Regorganizando términos ---que son escalares y se cumple la propiedad conmutativa
\[
  \vvv{u}\cdot\vvv{v}
  = \,u^ig_{ij}\,v^j
\]

Ahora, utilizamos la métrica para bajar el índice $j$, porque
$U_i = g_{ij}v^j$
\begin{equation}
  \vvv{u}\cdot\vvv{v}
  = u^i\,v_i = u_i\,v^i
\end{equation}


\section{Gradiente}
Supongamos un campo escalar $\phi$, definido en $\symbb{R^2}$ y expresado en coordenadas
oblicuas según sus componentes contravariantes ---utilizando, para la posición en el
espacio, la base unitaria $B= (\xhat{e}_1,\xhat{e}_2)$---
\begin{equation}
  \phi = \phi(x^1, x^2)
\end{equation}
Los vectores de esta base $B$ no cambian con la posición en el espacio.

\subsubsection{Componentes covariantes}
El gradiente es un vector covariante (1-forma). Sus componentes se obtienen derivando
parcialmente el campo escalar con respecto de las componentes contravariantes
$(\nabla\phi)_i = \partial\phi/\partial x^i$. Desarrollando la expresión anterior,
obtenemos
\begin{align}
  (\nabla\phi)_1 &= \dfrac{\partial\phi}{\partial x^1}\\
  (\nabla\phi)_2 &= \dfrac{\partial\phi}{\partial x^2}
\end{align}

El vector gradiente será
\begin{equation}
  \vvv{\nabla}\phi
  = (\nabla\phi)_i\vvv{e}^i
  = \dfrac{\partial\phi}{\partial x^i}\vvv{e}^i
  = \dfrac{\partial\phi}{\partial x^1}\vvv{e}^1
  + \dfrac{\partial\phi}{\partial x^2}\vvv{e}^2
\end{equation}

\subsubsection{Componentes contravariantes}
Para que el gradiente nos indique la dirección y sentido (segmento orientado) que indica
el máximo aumento del escalar tenemos que obtener sus componentes contravariantes.

Para obtener $\nabla\phi^1$ y $\nabla\phi^2$, tenemos que subir índices
$\nabla\phi_i=\partial\phi/\partial x^i$, multiplicando por la inversa de la métrica
\begin{equation}
  \vvv{\nabla}\phi
  = g^{ij}(\nabla\phi)_i\,\xhat{e}_j
  = g^{ij} \dfrac{\partial\phi}{\partial x^i}\,\xhat{e}_j
  = (\nabla\phi)^j\,\xhat{e}_j
\end{equation}

Desarrollando la expresión anterior
\begin{align*}
  \vvv{\nabla}\phi
  &=
    \mmm{g}^{11}\dfrac{\partial\phi}{\partial x^1}\,\xhat{e}_1
    + \mmm{g}^{12}\dfrac{\partial\phi}{\partial x^1}\,\xhat{e}_2
    + \mmm{g}^{21}\dfrac{\partial\phi}{\partial x^2}\,\xhat{e}_1
    + \mmm{g}^{22}\dfrac{\partial\phi}{\partial x^2}\,\xhat{e}_2\\
  &=
    \dfrac{1}{\sin^2\alpha}\,\dfrac{\partial\phi}{\partial x^1}\,\xhat{e}_1
    - \dfrac{\cos\alpha}{\sin^2\alpha}\,\dfrac{\partial\phi}{\partial x^1}\,\xhat{e}_2
    - \dfrac{\cos\alpha}{\sin^2\alpha}\,\dfrac{\partial\phi}{\partial x^2}\,\xhat{e}_1
    + \dfrac{1}{\sin^2\alpha}\,\dfrac{\partial\phi}{\partial x^2}\,\xhat{e}_2
\end{align*}

Simplificando, obtendríamos
\begin{equation}\label{eq:R2-oblic_gradiente_contrav}
  \vvv{\nabla}\phi
  = \dfrac{1}{\sin^2\alpha}
  \left[%
    \left(%
      \dfrac{\partial\phi}{\partial x^1}-\cos\alpha\dfrac{\partial\phi}{\partial x^2}
    \right) \xhat{e}_1
    +\left(%
      \dfrac{\partial\phi}{\partial x^2}-\cos\alpha\dfrac{\partial\phi}{\partial x^1}
    \right) \xhat{e}_2
  \right]
\end{equation}

{\small
  También se podrían haber calculado subiendo el índice de las componentes por separado
\begin{equation}
  (\nabla\phi)^i
  = g^{ij} (\nabla\phi)_j = g^{ij}\dfrac{\partial\phi}{\partial x^j}
\end{equation}
}

{\small
\begin{align*}
  (\nabla\phi)^1
  &=
    g^{11} \dfrac{\partial\phi}{\partial x^1}
    + g^{12} \dfrac{\partial\phi}{\partial x^2}
    =
    \dfrac{1}{\sin^2\alpha}\dfrac{\partial\phi}{\partial x^1}
    - \dfrac{\cos\alpha}{\sin^2\alpha}\dfrac{\partial\phi}{\partial x^2}
    =
    \dfrac{1}{\sin^2\alpha}
    \left(%
    \dfrac{\partial\phi}{\partial x^1}-\cos\alpha\dfrac{\partial\phi}{\partial x^2}
    \right)\\
  (\nabla\phi)^2
  &=
    g^{21} \dfrac{\partial\phi}{\partial x^1}
    + g^{22} \dfrac{\partial\phi}{\partial x^2}
    =
    - \dfrac{\cos\alpha}{\sin^2\alpha}\dfrac{\partial\phi}{\partial x^1}    
    \dfrac{1}{\sin^2\alpha}\dfrac{\partial\phi}{\partial x^2}
    =
    \dfrac{1}{\sin^2\alpha}
    \left(%
    \dfrac{\partial\phi}{\partial x^2}-\cos\alpha\dfrac{\partial\phi}{\partial x^1}
    \right)  
\end{align*}
}

{\small
Y luego sustituir lo anterior en
\begin{equation}
  \vvv{\nabla\phi} = (\nabla\phi)^1\,\xhat{e}_1 + (\nabla\phi)^2\,\xhat{e}_2
\end{equation}
dando el mismo resultado que en \eqref{eq:R2-oblic_gradiente_contrav}.
}






%%% Local Variables:
%%% mode: latex
%%% TeX-engine: luatex
%%% TeX-master: "../retazosmatematicas.tex"
%%% End:
